\providecommand{\toplevelprefix}{../..}  %
\documentclass[../../book-main.tex]{subfiles}

\begin{document}

\chapter{Las representaciones profundas como optimización desenrollada}
\label{ch:deep-networks}

\begin{quote}
    \hfill    ``{\em Lo que no puedo crear, no lo entiendo}.''

    $~$ \hfill --- Richard Feynman
\end{quote}
\vspace{5mm}

En capítulos anteriores, hemos mostrado cómo identificar estructuras de baja dimensión
en espacios de alta dimensión, \textit{centrándonos principalmente en
estructuras lineales}.
Por ejemplo, introdujimos el análisis de componentes principales (PCA) para
aprender el filtro de ruido lineal $\hat{\bm{U}}^{\top}$ cuando los datos observados
$\bm{x}$ siguen el modelo estadístico $\bm{x} = \bm{U}\bm{z} +
\bm{\varepsilon}$.
En este contexto, las representaciones aprendidas los datos de entrada linealmente
transformados $\hat{\bm{U}}^{\top}\bm{x}$.
Bajo el supuesto de un modelo lineal, es posible aprender la estructura lineal
de baja dimensión mediante algoritmos de optimización eficientes y con sólidas
garantías teóricas.
Más aún, el supuesto del modelo lineal abarca una amplia gama de
aplicaciones y problemas, entre los que se incluyen el reconocimiento facial, la
recuperación de imágenes por resonancia magnética y la recuperación de texturas estructurales~\cite{Wright-Ma-2022}.

Por otro lado, el modelo lineal puede resultar limitado cuando se trata de
aplicaciones del mundo real, especialmente cuando los datos de entrada $\bm{x}$ son
complejos, como el habla y los lenguajes naturales, las imágenes y los videos, así como
los movimientos robóticos. Las distribuciones de baja dimensión de dichos datos suelen ser
{\em no lineales}.
El tratamiento de la no linealidad cuenta con una larga trayectoria en diversas
disciplinas, como la teoría de control, el procesamiento de señales y el reconocimiento de patrones.
Se han realizado esfuerzos considerables para tratar de extender
los métodos y soluciones de los modelos lineales con el fin de abordar la no linealidad,
incluidos los primeros intentos de ampliar el PCA al PCA no lineal (como veremos
con más detalle en el \Cref{ch:consistent-representations}). En
la mayoría de los casos, los métodos se diseñan partiendo de 
ciertas suposiciones sobre la distribución de los datos
y se adaptan a problemas específicos.

Más recientemente, las redes neuronales profundas han logrado un éxito notable
en una amplia variedad de datos y aplicaciones. Una red neuronal
\begin{equation}
    f(\cdot, \bm \theta) \colon \x
    \xrightarrow{\hspace{1mm} f^{\mathrm{pre}} \hspace{1mm}} \z^1
    \rightarrow \cdots
    \rightarrow \z^\ell \xrightarrow{\hspace{1mm} f^\ell \hspace{1mm}}
    \z^{\ell+1} \rightarrow  \cdots \to \z^{L+1} = \z.
\end{equation}
puede aprender características o representaciones efectivas para aplicaciones posteriores.
Por ejemplo, una red neuronal profunda entrenada $f(\cdot,
\bm \theta)$ puede aplicarse para mapear imágenes a vectores de características, es
decir, $\bm{z}_i = f(\bm{x}_i,\bm \theta)$, mientras que un clasificador lineal
puede aprenderse sobre tales representaciones $\{\bm{z}_i\}$. Un
avance notable es AlexNet~\cite{krizhevsky2012imagenet}, una red neuronal
convolucional profunda entrenada con más de un millón de imágenes
naturales, que supera (en tareas de predicción) a todos los enfoques
anteriores basados en características definidas manualmente. Una de las diferencias clave
entre AlexNet y los enfoques anteriores es que el
primero \textit{aprende los parámetros de la transformación no lineal a partir de
enormes cantidades de datos} entrenados mediante retropropagación (BP)
\cite{Back-Prop}, tal y como se detalla en el \Cref{app:BP-section} del
\Cref{app:optimization}.

En la práctica habitual posterior, se modela el mapeo $f$ con otras
redes neuronales artificiales profundas diseñadas empíricamente y se aprenden los
parámetros $\bm \theta$ a partir de una inicialización aleatoria mediante el algoritmo 
de retropropagación. Desde
la aparición de AlexNet \cite{krizhevsky2012imagenet}, las arquitecturas de
las redes profundas modernas siguen siendo revisadas y mejoradas de forma empírica.
Arquitecturas de red como VGG \cite{simonyan2014very}, ResNet
\cite{he2016deep}, DenseNet \cite{dense-net}, CNN, RNN o LSTM
\cite{LSTM}, Transformer \cite{vaswani2017attention} y mezclas
de expertos (MoE) \cite{MoE,Fedus-2022}, etc., han seguido ampliando
los límites del rendimiento. Como parte del esfuerzo por mejorar el
rendimiento de las redes profundas, se ha analizado empíricamente casi todos los componentes de las redes
y se han propuesto diversas revisiones y mejoras. No se limitan a funciones de activación no lineales
\cite{maas2013rectifier,klambauer2017self,xu2015empirical,nwankpa2018activation},
conexiones de salto \cite{ronneberger2015u,he2016deep}, normalizaciones
\cite{ioffe2015batch,ba2016layer,ulyanov2016instance,wu2018group,miyato2018spectral},
muestreo ascendente/descendente o agrupación \cite{scherer2010evaluation},
convoluciones~\cite{lecun1998gradient,krizhevsky2012imagenet}, etc. Sin embargo, casi 
todas estas modificaciones se han desarrollado a lo largo de
años de estudios empíricos {de prueba y error} o de ablación. Algunas prácticas recientes
llegaron incluso al extremo de buscar estructuras de red
y estrategias de entrenamiento efectivas mediante técnicas de búsqueda aleatoria
extensiva, como la Búsqueda de Arquitectura Neuronal
\cite{NAS-1,Baker2017DesigningNN}, AutoML \cite{automl} y Aprendizaje
para Aprender \cite{andrychowicz2016learning}.

A pesar de la amplia aplicación de las redes neuronales profundas, no está claro
cuáles son los principios de diseño subyacentes de una red construida de este modo.
En particular, no está claro qué función matemática desempeña cada
capa de la red. En este capítulo, basándonos en los resultados
de los capítulos anteriores, desarrollamos un marco teórico que
proporcionará una interpretación matemática totalmente rigurosa del papel de una
red profunda, incluyendo sus capas individuales y la red en su conjunto.

Para comprender las redes profundas y cómo se pueden diseñar mejor,
debemos partir del objetivo del aprendizaje de representaciones. En
los capítulos anteriores, hemos argumentado que el objetivo consiste en identificar
la distribución de los datos, que es intrínsecamente de baja dimensión, y luego
transformarla en una representación compacta y estructurada (por ejemplo, lineal por tramos).
Como hemos visto en el capítulo anterior, el enfoque general
para identificar una distribución de datos de baja dimensión consiste
en un proceso de compresión que minimiza progresivamente la
entropía o la tasa de codificación de la distribución. Sin embargo, hasta ahora
hemos estado utilizando redes profundas diseñadas de forma empírica para modelar
o aproximar las operaciones destinadas a optimizar estos objetivos,
como la función de puntuación para la eliminación de ruido (en
la \Cref{sec:denoising-intro}) o la transformación que maximiza la
reducción de la tasa (en la \Cref{sec:MCR-landscape}).

Como hemos argumentado en el capítulo anterior,
en la \Cref{sec:chap4-representation-learning-problem} en particular, se
puede evaluar la calidad de la representación resultante mediante la
información que esta aporta con respecto a una representación «perezosa»
que modela todos los datos como una gran
distribución gaussiana.\footnote{que hemos visto en el capítulo anterior como una
opción concreta de interpretación del conjunto de datos muestreado.} En
particular, si \textit{utilizamos una mezcla de gaussianas
    (subespacios)\footnote{que hemos estudiado a fondo en el capítulo
    anterior} como distribuciones prototípicas para aproximar la distribución no lineal
que nos interesa}, entonces podemos medir de manera eficiente la tasa de codificación
de dicha representación utilizando la (suma de las) funciones de tasa-distorsión
de las gaussianas asociadas. Entonces, la cantidad de {\em
información ganada} o la entropía (relativa) reducida con dicho
modelado puede medirse mediante la diferencia entre la tasa de codificación
de la representación simplificada y la de la representación más refinada.
Por lo tanto, el objetivo del aprendizaje de representaciones es
maximizar esta ganancia de información, también conocida como el objetivo de reducción de tasa.

Como veremos en este capítulo, una vez que el objetivo del aprendizaje de representaciones
queda claro, la función de una red neuronal profunda consiste precisamente en
ayudar a optimizar ese objetivo de forma iterativa. Cada capa de una red neuronal profunda
puede derivarse de forma natural como un paso de optimización iterativo destinado a
maximizar de manera incremental la ganancia de información, incluyendo las populares
arquitecturas de ResNet, CNN y Transformadores, así como otras variantes más
avanzadas. En particular, este capítulo tiene como objetivo responder a las
siguientes preguntas sobre las redes profundas:
\begin{itemize}
    \item \Cref{sec:chap4-white-box-model-via-unrolling} --- dada una
        medida de calidad para una representación aprendida, ¿cómo
        construir el mapeo no lineal entre los datos y la representación óptima
        mediante la optimización desenrollada del objetivo?
    \item \Cref{sec:chap4-white-box-transformer} --- ¿cómo podría el
        enfoque de «desenrollado» mencionado anteriormente ofrecer una interpretación fundamentada
        de las populares arquitecturas de transformadores? Si es así, ¿cuáles son los
        mecanismos de optimización y los objetivos asociados?
    \item \Cref{sec:chap4-derive-white-box-transformer-variants} ---
        ¿cómo nos ayudaría este marco a diseñar arquitecturas profundas más eficientes o
        más parsimoniosas?
\end{itemize}

\section{Redes profundas de white-box mediante optimización
desenrollada}\label{sec:chap4-white-box-model-via-unrolling}

Ahora bien, si estamos de acuerdo en que maximizar la reducción de la tasa o la ganancia de información
conduce a la representación deseada, tal como se analiza en la
\Cref{sec:chap4-representation-learning-problem}, la pregunta que queda por resolver
es cómo construir y aprender un mapeo (no lineal) entre los
datos $\X$ y la representación óptima $\Z^*$. Esto implica
diseñar una arquitectura de red y un algoritmo de aprendizaje que puedan
captar de manera eficaz las estructuras subyacentes de los datos y
reproducir fielmente la representación óptima.

\subsection{Redes profundas a partir del descenso por gradiente desenrollada}

En el capítulo anterior, presentamos el objetivo de reducción de la tasa
\eqref{eqn:maximal-rate-reduction} como un objetivo fundamentado para
aprender representaciones discriminativas lineales de los datos. Sin embargo,
no hemos especificado la arquitectura de la función de mapeo de características $\z =
f(\x, \bm \theta)$ para extraer dichas representaciones de los datos de entrada $\x$.
Una opción sencilla es utilizar una red profunda convencional, como
ResNet, para implementar $f(\x, \bm \theta)$. Como hemos visto en el
\Cref{eg:Rate-Reduction-CIFAR10}, esta opción suele dar lugar empíricamente a un rendimiento
aceptable en la práctica. No obstante, siguen existiendo varios problemas sin resolver
al adoptar una red profunda arbitraria. Aunque la
representación de características aprendida es ahora más interpretable, la red
en sí misma sigue {\em sin} serlo. No está claro por qué cualquier red de «black-box»
elegida es capaz de optimizar el objetivo MCR$^2$ deseado. Los
buenos resultados empíricos (por ejemplo, con una ResNet) no justifican necesariamente
la elección concreta de arquitecturas y operadores de la red: 
¿Por qué es necesario un modelo de capas profundas? ¿Qué es lo que las capas adicionales
pretenden mejorar o simplificar? ¿Qué anchura y profundidad son adecuadas? ¿O existe
alguna justificación rigurosa para las convoluciones (en su popular
forma multicanal) y los operadores no lineales (por ejemplo, ReLU o softmax) que se utilizan?

En este capítulo, demostramos que el uso del ascenso por gradiente para maximizar la
reducción de la tasa $\Delta R_{\epsilon}(\Z \mid \bm \Pi)$, tal como se define en
\eqref{eqn:maximal-rate-reduction}, conduce de manera natural a una red profunda de «white-box»
que realiza el mapeo deseado. Todas las capas de la red, los operadores lineales y no lineales, 
así como los parámetros se construyen {\em explícitamente
siguiendo un enfoque de propagación puramente hacia adelante}. Además, estas
arquitecturas de red se asemejan a las redes profundas existentes diseñadas
empíricamente, lo que proporciona justificaciones fundamentadas para su diseño.

\paragraph{Ascenso por gradiente para la reducción de la tasa de codificación.} En el
capítulo anterior vimos que, para buscar una representación discriminativa lineal
(LDR), matemáticamente estamos buscando, en esencia, una
función continua $f(\cdot): \x \mapsto \z$ que mapee los datos $\X =
[\x_1, \ldots, \x_N] \in \Re^{D \times N}$ (o características iniciales
    extraídas de los datos\footnote{Como veremos, la necesidad de dicha
extracción de características se abordará en la siguiente sección.}) a una representación óptima
$\Z = [\z_1, \ldots, \z_N] \in \Re^{d \times N}$ que
maximice el siguiente objetivo de reducción de la tasa de codificación (véase también
\eqref{eq:MCR pi}):
\begin{equation}\label{eq:mcr-formulation}
    \scalebox{0.95}{\(\displaystyle\Delta R_{\epsilon}(\Z \mid \bm{\Pi}) \doteq
            \underbrace{\frac{1}{2}\log\det \Big(\I + {\alpha} \Z
            \Z^{\top} \Big)}_{R_{\epsilon}(\Z)} \;-\;
            \underbrace{\sum_{k=1}^{K}\frac{\gamma_{k}}{2}
                \log\det\Big(\I + {\alpha_{k}} \Z \bm{\Pi}_{k} \Z^{\top}
    \Big)}_{R_{\epsilon}^c(\Z \mid \bm \Pi)}\)},
\end{equation}
donde $\epsilon > 0$ es un error de cuantificación predeterminado y, para
simplificar, denotamos\footnote{Obsérvese que utilizamos una notación ligeramente simplificada
en comparación con el \Cref{ch:denoising}.}
\begin{equation}
    \alpha \doteq \frac{d}{N\epsilon^2}, \quad \alpha_{k} \doteq
    \frac{d}{\mathrm{tr}(\bm{\Pi}_{k})\epsilon^2}, \quad \gamma_{k}
    \doteq \frac{\mathrm{tr}(\bm{\Pi}_{k})}{N}, \quad \text{for}\ k
    = 1,\ldots, K.
\end{equation}

La pregunta se reduce, en realidad, a si existe una forma {\em
constructiva} de encontrar una aplicación continua $f(\cdot,\bm
\theta)$ de $\bm x$ a $\bm z$. Para ello, consideremos
la maximización incremental del objetivo $\Delta R_{\epsilon}$ como una
función de $\Z \subseteq \mathbb{S}^{d-1}$. Aunque puede haber
muchos esquemas de optimización entre los que elegir, por simplicidad consideraremos primero
el esquema de «ascenso por gradiente proyectado» (PGA), posiblemente el más sencillo:\footnote{Obsérvese que utilizamos el subíndice $j$ en $\Z_j$ para
    indicar las características de la $j$-ésima clase y el superíndice $\ell$ en
$\Z^{\ell}$ para indicar todas las características en la iteración o capa $\ell$-ésima.}
\begin{equation}
    \bm Z^{\ell+1}   \; \propto \; \bm Z^{\ell} + \eta \cdot
    \frac{\partial \Delta R_{\epsilon}}{\partial \bm Z}(\Z^{\ell})
    \quad \mbox{s.t.} \quad \Z^{\ell+1} \subseteq \mathbb{S}^{d-1},
    \quad \ell = 1, 2, \ldots,
    \label{eqn:gradient-descent}
\end{equation}
para un tamaño de paso $\eta > 0$ y el iterado inicial dado
por $\bm Z^{0} = \bm X$.\footnote{Una vez más, para simplificar, aquí
    suponemos en primer lugar que las características iniciales $\bm Z^{1}$ son los propios datos.
    Obsérvese que aquí $\ell$ denota el número de iteraciones.
    Por lo tanto, los datos y las características tienen la misma dimensión $d$. Sin embargo,
    esto no tiene por qué ser así. Como veremos en la siguiente sección,
    las características iniciales pueden ser algunas características (elevadas) de los datos para
    empezar y, en principio, podrían tener una dimensión diferente (mucho mayor). 
Todas las iteraciones posteriores tienen la misma dimensión.}
Este esquema puede interpretarse como la forma en que se deben ajustar gradualmente
las ubicaciones de las características actuales $\Z^\ell$, inicializadas a partir de los datos de entrada
$\bm X$, para que el conjunto resultante $\Z^{\ell +1}$ mejore
la tasa de reducción $\Delta R_{\epsilon}$, tal como se ilustra en la
\Cref{fig:gradient-flow}.
\begin{figure}
    \centering
    \includegraphics[width=0.85\linewidth]{\toplevelprefix/chapters/chapter5/figs/redu_gradient_diagram.png}
    \caption{Deformación incremental mediante flujo de gradiente para, por un lado,
        aplanar los datos de cada clase en un subespacio y, por otro,
    separar las diferentes clases.}
    \label{fig:gradient-flow}
\end{figure}

Un cálculo sencillo muestra que el gradiente ${\partial \Delta
R_{\epsilon}}/{\partial \bm Z}$ implica evaluar las siguientes
derivadas de los dos términos de $\Delta R_{\epsilon}$:
\begin{equation}\label{eqn:expand-directions}
    \frac{1}{2}\frac{\partial \log \det (\I \!+\! \alpha \Z \Z^{\top}
    )}{\partial \bm Z}(\Z^\ell) = \underbrace{\alpha(\I \!+\!
    \alpha\Z^\ell (\Z^\ell)^\top)^{-1}}_{\bm E^{\ell} \; \in
    \Re^{d\times d}}\Z^\ell,
\end{equation}

\begin{equation}\label{eqn:compress-directions}
    \scalebox{0.925}{\(\displaystyle\frac{1}{2}\frac{\partial \left(
                    \gamma_{k}  \log \det (\I +
            \alpha_{k} \Z \bm \Pi_k \Z^{\top} )  \right)}{\partial \bm
            Z}(\Z^\ell) = \gamma_{k}  \underbrace{ \alpha_{k}  (\I +
            \alpha_{k} \Z^\ell \bm \Pi_k (\Z^\ell)^{\top})^{-1}}_{\bm
    C^{\ell}_k \; \in \Re^{d\times d}} \Z^{\ell} \bm \Pi_k\)}.
\end{equation}
Nótese que, en lo anterior, la matriz $\bm E^{\ell}$ solo depende de
$\Z^{\ell}$ y tiene como objetivo {\em expandir} todas las características para aumentar
la tasa de codificación global; la matriz $\bm C^{\ell}_{k}$ depende de
las características de la clase $k$ y tiene como objetivo {\em comprimirlas} para reducir
la tasa de codificación de cada clase.
Entonces, el gradiente completo $\frac{\partial \Delta
R_{\epsilon}}{\partial \bm Z}(\Z^\ell) \in \Re^{d\times N}$ tiene la  forma:
\begin{equation}
    \frac{\partial \Delta R_{\epsilon}}{\partial \bm Z}(\Z^\ell)  =
    \underbrace{\bm E^{\ell}}_{\text{Expansion}} \Z^{\ell} \;-\;
    \sum_{k=1}^K \gamma_{k} \underbrace{\bm
    C_{k}^{\ell}}_{\text{Compression}}  \Z^{\ell} \bm{\Pi}_k.
    \label{eqn:DR-gradient}
\end{equation}

\begin{remark}[Interpretación de $\bm E^\ell$ y $\bm C_j^\ell$ como
    operadores lineales]\label{rem:regression-interpretation}
    Para cualquier $\bm z^\ell \in \mathbb{R}^d$,
    \begin{gather}
        \scalebox{0.95}{\(\displaystyle\bm E^\ell \bm z^\ell =
                \alpha(\bm z^\ell - \bm Z^\ell \bm
                q^\ell_\star), \quad
                \mbox{donde} \quad \bm q^\ell_\star \doteq \argmin_{\bm
                q^\ell} \big\{\alpha \|\bm z^\ell - \bm Z^\ell \bm
        q^\ell\|_2^2 + \|\bm q^\ell\|_2^2\big\}.\)}
    \end{gather}
    Obsérvese que $\bm q^\ell_\star$ es exactamente la solución de la
    regresión de crestas para todos los puntos de datos $\bm Z^\ell$ en cuestión.
    Por lo tanto, $\bm E^\ell$ (y de manera similar para $\bm C^\ell_k$) es
    aproximadamente (es decir, cuando $N$ es lo suficientemente grande) la proyección
    sobre el complemento ortogonal del subespacio generado por las columnas
    de $\bm Z^\ell$. Otra forma de interpretar la matriz $\bm E^\ell$
    es mediante la descomposición en valores propios de la matriz de covarianza
    $\Z^\ell (\Z^\ell)^\top$. Suponiendo que $\Z^\ell (\Z^\ell)^\top
    \doteq \bm U^\ell \bm \Lambda^\ell (\bm U^\ell)^\top$, donde $\bm
    \Lambda^\ell \doteq \diag\left(\lambda^\ell_1, \ldots,
    \lambda^\ell_d \right)$, tenemos
    \begin{equation}
        \bm E^\ell = \alpha \bm U^\ell\,
        \diag\left(\frac{1}{1+\alpha\lambda^\ell_1}, \ldots,
        \frac{1}{1+\alpha\lambda^\ell_d}\right) \left(\bm U^\ell\right)^\top.
    \end{equation}
    Por lo tanto, la matriz $\bm E^\ell$ actúa sobre un vector $\bm
    z^\ell$ estirándolo de tal manera que las direcciones de gran varianza
    se reducen, mientras que se mantienen las direcciones de varianza cercana a cero. Estas
    son precisamente las direcciones \eqref{eqn:expand-directions} en las que
    desplazamos las características para que el volumen total se expanda y la
    tasa de codificación aumente, de ahí el signo positivo. Por el
    contrario, las direcciones asociadas a
    \eqref{eqn:compress-directions} son «residuos» de las características de
    cada clase que se desvían del subespacio al que se
    supone que pertenecen. Estas son exactamente las direcciones en las que las
    características deben comprimirse de nuevo hacia su respectivo
    subespacio, de ahí el signo negativo (véase la
    \Cref{fig:regression-interpretation}).

    \begin{figure}[t]
        \centering
        \includegraphics[width=0.65\linewidth]{\toplevelprefix/chapters/chapter5/figs/expand_compress.png}
        \caption{\small Interpretación de $\bm C^\ell_k$ y $\bm
            E^\ell$: $\bm C^\ell_k$ comprime cada clase contrayendo
            las características a un subespacio de baja dimensión; $\bm E^\ell$
            expande todas las características contrastando y repeliendo características
        entre diferentes clases.}
        \label{fig:regression-interpretation}
    \end{figure}

    En esencia, las operaciones lineales $\bm E^\ell$ y $\bm
    C_k^\ell$ en el método de ascenso por gradiente para la reducción de la tasa se determinan mediante
    los datos de entrenamiento que realizan «autorregresiones». Los recientes avances
    en la comprensión de la regresión de crestas en un contexto sobreparametrizado
    \cite{yang2020rethinking,Wu2020OnTO} indican que el uso
    de datos aparentemente muestreados de forma redundante (de cada subespacio) como
    regresores no conduce a sobreajuste.
\end{remark}

\paragraph{Incremento del mapa de características guiado por gradientes.} Nótese que en lo
anterior, el ascenso por gradiente considera todas las características $\Z^{\ell} =
[\z^{\ell}_{1}, \dots, \z^{\ell}_{N}]$ como variables libres. El
incremento $\Z^{\ell+1} - \Z^{\ell} = \eta \frac{\partial \Delta
R_{\epsilon}}{\partial \bm Z}(\Z^\ell)$ aún no proporciona una
transformación en todo el dominio de características $\z^\ell \in \Re^d$.
Según la ecuación \eqref{eqn:DR-gradient}, el gradiente no puede
calcularse en un punto cuya pertenencia se desconoce, tal como se ilustra en la
\Cref{fig:gradient-flow}. Por lo tanto, para hallar la
$f(\x,\bm  \theta)$ de forma explícita, podemos considerar la construcción de una
transformación de incremento pequeño $g(\cdot, \theta^{\ell})$ sobre la característica de la
$\ell$-ésima capa $\z^\ell$ para emular el esquema de gradiente (proyectado) anterior:
\begin{equation}
    \z^{\ell + 1}   \; \propto \; \z^{\ell} + \eta\cdot  g(\z^{\ell},
    \theta^{\ell}) \quad \mbox{sujeto a} \quad \z^{\ell +1}
    \in \mathbb{S}^{d-1}
    \label{eqn:gradient-descent-transform}
\end{equation}
de tal manera que $\big[g(\z_1^{\ell}, \bm \theta^{\ell}), \ldots,
g(\z_N^{\ell}, \bm \theta^{\ell}) \big] \approx \frac{\partial \Delta
R_{\epsilon}}{\partial \bm Z}(\Z^\ell).$ Es decir, necesitamos
aproximar el flujo del gradiente $\frac{\partial \Delta
R_{\epsilon}}{\partial \bm Z}$ que deforma localmente todas las características (de entrenamiento)
$\{\z_i^\ell\}_{i=1}^N$ con un mapeo continua $g(\z,\bm
\theta)$ definida en todo el espacio de características $\z^\ell \in \Re^d$.
Obsérvese que el incremento
\eqref{eqn:gradient-descent-transform} puede interpretarse como una versión discretizada de una
ecuación diferencial continua:
\begin{equation}
    \dot{\z} = g(\z, \theta).
\end{equation}
Por lo tanto, la red (profunda) así construida puede interpretarse como una
determinada ecuación diferencial ordinaria neuronal \cite{chen2018neural}. Sin embargo, a diferencia de la
ODE neuronal, donde el flujo $g$ se elige para que sea una estructura genérica, aquí
nuestro $g(\z, \bm\theta)$ emula el flujo de gradiente de la reducción de la tasa
en el conjunto de características (como se muestra en la \Cref{fig:gradient-flow}):
\begin{equation*}
    \dot{\Z} = \frac{\partial \Delta R_{\epsilon}}{\partial \bm Z},
\end{equation*}
y su estructura se deriva y se determina íntegramente a partir de este
objetivo, sin ningún otro supuesto previo ni heurística.

Al examinar la estructura del gradiente \eqref{eqn:DR-gradient},
se deduce que una opción natural para la transformación de incremento
$g(\z^\ell, \bm \theta^\ell)$ es de la forma:
\begin{equation}
    g(\z^\ell, \bm \theta^\ell) \; \doteq \; \bm E^\ell \z^\ell -
    \sum_{k=1}^K \gamma_{k} \pi_k(\z^\ell)\bm C_{k}^{\ell}  \z^{\ell}
    \in \Re^d,
    \label{eqn:DR-gradient-transform}
\end{equation}
donde $\pi_k(\z^\ell) \in [0,1]$ indica la probabilidad de que
$\z^{\ell}$ pertenezca a la clase $k$-ésima. Los parámetros del mapa de incremento
$\bm \theta^\ell$ dependen: En primer lugar, de un conjunto de mapas lineales
representados por $\bm E^{\ell}$ y $\{ \bm C^{\ell}_{k}\}_{k=1}^{K}$
que dependen únicamente de las estadísticas de las características del entrenamiento $\Z^\ell$.
En segundo lugar, de la pertenencia $\{ \pi_k(\z^\ell)\}_{k=1}^K$ de cualquier característica $\z^\ell$.
Nótese que, en las muestras de entrenamiento $\Z^\ell$, para las que se conocen las
pertenencias $\bm \Pi_k$, la función $g(\z^\ell, \bm
\theta)$ así definida proporciona exactamente los valores del gradiente $\frac{\partial
\Delta R_{\epsilon}}{\partial \bm Z}(\Z^\ell)$.

Dado que solo disponemos de la función de pertenencia para los datos de entrenamiento, la
función $g(\cdot)$ definida en \eqref{eqn:DR-gradient-transform} solo puede
calcularse sobre los datos de entrenamiento. Para extrapolar $g(\cdot)$ a todo el
espacio de características, necesitamos estimar $\pi_k(\z^\ell)$ en su
segundo término. En el aprendizaje profundo convencional, este mapa suele
modelarse como una red profunda y aprenderse a partir de los datos de entrenamiento, por ejemplo, mediante
{\em la retropropagación}. Sin embargo, nuestro objetivo aquí no es aprender ya un
clasificador preciso $\pi_{k}(\z^\ell)$. En cambio, solo necesitamos
una estimación lo suficientemente buena de la información de clase para que
$g(\cdot)$ aproxime bien el gradiente $\frac{\partial \Delta
R_{\epsilon}}{\partial \bm Z}$.

A partir de la interpretación geométrica de las aplicaciones lineales $\bm E^\ell$ y
$\bm C_k^\ell$ dada por la \Cref{rem:regression-interpretation}, el
término $\bm p_{k}^{\ell} \doteq \bm C^{\ell}_k \z^{\ell}$ puede considerarse
como (aproximadamente) la proyección de $\z^{\ell}$ sobre el complemento ortogonal
de cada clase $j$. Por lo tanto, $\|\bm p_{j}^{\ell}\|_2$ es
pequeño si $\z^\ell$ está en la clase $j$ y grande en caso contrario. Esto
nos lleva a estimar su pertenencia basándonos en la siguiente función softmax:
\begin{align}
    \widehat{\vpi}(\z^\ell) &\doteq
    \softmax\rp{-\lambda\mat{\norm{\vC^{\ell}_{1}\vz^{\ell}}_{2}
    \\ \vdots \\ \norm{\vC^{\ell}_{K}\vz^{\ell}}_{2}}} \\
    &= \frac{1}{\sum_{k = 1}^{K}\exp(-\lambda
    \norm{\vC^{\ell}_{k}\vz^{\ell}}_{2})}\mat{\exp(-\lambda\norm{\vC^{\ell}_{1}\vz^{\ell}}_{2})
    \\ \vdots \\ \exp(-\lambda\norm{\vC^{\ell}_{K}\vz^{\ell}}_{2})}
    \in [0,1]^{K}.
\end{align}
Por lo tanto, el segundo término de \eqref{eqn:DR-gradient-transform} puede
aproximarse mediante esta pertenencia estimada:
\begin{equation}
    \sum_{k=1}^K \gamma_{k} \pi_k(\z^\ell)\bm C_{k}^{\ell}  \z^{\ell}
    \; \approx \;  \sum_{k=1}^K \gamma_{k} \widehat{\pi}_k(\z^\ell)
    \bm{C}^{\ell}_k  \z^{\ell}
    \; \doteq \; \bm \sigma\Big([\bm{C}^{\ell}_{1} \z^{\ell}, \dots,
    \bm{C}^{\ell}_{K} \z^{\ell}]\Big),
    \label{eqn:soft-residual}
\end{equation}
que se denota como un operador no lineal $\bm \sigma(\cdot)$ sobre
las salidas de la característica $\z^\ell$ a través de $K$ grupos de filtros:
$[\bm{C}^{\ell}_{1}, \dots, \bm{C}^{\ell}_{K}]$. Obsérvese que la
no linealidad surge debido a una asignación «suave» de la pertenencia a una clase
basada en las respuestas de las características de esos filtros.

En general, al combinar \eqref{eqn:gradient-descent-transform},
\eqref{eqn:DR-gradient-transform} y \eqref{eqn:soft-residual},
la transformación incremental de características de $\z^{\ell}$ a $\z^{\ell+1}$ queda ahora así:
\begin{align}\label{eqn:layer-approximate}
    \z^{\ell+1}  &\propto \; \z^\ell +  \eta \cdot  \bm E^{\ell}
    \z^{\ell} - \eta\cdot  \bm \sigma\big([\bm{C}^{\ell}_{1}
    \z^{\ell}, \dots, \bm{C}^{\ell}_{K} \z^{\ell}]\big)  \\
    &= \; \z^\ell +  \eta \cdot g(\z^\ell, \bm \theta^\ell)
    \qquad \mbox{s.t.} \quad \z^{\ell +1} \in \mathbb{S}^{d-1},
\end{align}
con la función no lineal $\bm \sigma(\cdot)$ definida anteriormente y
$\bm \theta^\ell$ que recoge todos los parámetros por capa. Es decir,
$\bm \theta^\ell =\left\{\bm E^\ell, \bm{C}^{\ell}_{1}, \dots,
\bm{C}^{\ell}_{K}, \gamma_{k}, \lambda\right\}$. Tenga en cuenta que las características de
cada capa siempre se «normalizan» mediante la proyección sobre la esfera
unitaria $\mathbb S^{d-1}$, denotada como $\mathcal P_{\mathbb S^{d-1}}$,
es decir, dividiendo cada característica por su norma.\footnote{Esto refleja el
    bloque LayerNorm habitual en las redes neuronales, que calcula la
    característica proyectada y, a continuación, la multiplica y suma mediante parámetros
    entrenables \citep{ba2016layer}; véase la
\Cref{sub:contrastive_learning_architecture} para obtener más detalles.}
La forma del incremento en \eqref{eqn:layer-approximate} puede
ilustrarse mediante un diagrama en la \Cref{fig:arch-ReduNet}(a).

\begin{figure}[t]
    \begin{subfigure}[t]{0.4\textwidth}
        \centering
        \includegraphics[width=\textwidth]{\toplevelprefix/chapters/chapter5/figs/redunet_layer.png}
        \caption{\textbf{ReduNet}}
    \end{subfigure}
    \hfill
    \begin{subfigure}[t]{0.58\textwidth}
        \centering
        \includegraphics[width=\textwidth]{\toplevelprefix/chapters/chapter5/figs/resnet_resnext.pdf}
        \caption{\textbf{ResNet} y \textbf{ResNeXt}.}
    \end{subfigure}
    \caption{\small Arquitecturas de red de ReduNet y
        comparación con otras. \textbf{(a)}: Estructura de capas de
        \textbf{ReduNet} derivada de una iteración de ascenso por gradiente
        para optimizar la reducción de la tasa. \textbf{(b)} (izquierda): Una capa de
        ResNet~\cite{he2016deep}; y \textbf{(b)} (derecha): Una capa de
        ResNeXt~\cite{ResNEXT}. Como veremos en la
        \Cref{sec:shift-invariant}, los operadores lineales $\bm E^\ell$ y
        $\bm{C}_k^\ell$ de ReduNet se convierten naturalmente en convoluciones (multicanal)
    cuando se impone la invariancia al desplazamiento.}
    \label{fig:arch-ReduNet}
\end{figure}

\begin{algorithm}[t]
    \caption{Algoritmo de entrenamiento para ReduNet}
    \label{alg:redunet_training}
    \begin{algorithmic}[1]
        \Require{$\vX = [\vx_1, \ldots, \vx_N]\in \Re^{D \times N}$,
            $\bm{\Pi} = \{\bm \Pi_k\}_{k=1}^K$, $\epsilon > 0$,
        $\lambda$ y una tasa de aprendizaje $\eta$.}
        \Ensure{Los parámetros aprendidos \(\{\vE^{\ell}\}_{\ell =
                1}^{L}, \{\{\vC^{\ell}_{k}\}_{k = 1}^{K}\}_{\ell = 1}^{L},
        \{\gamma_{k}\}_{k = 1}^{k}\)}

        \Procedure{ReduNetTraining}{$\vX, \vPi, \epsilon, \lambda, \eta$}
        \Statex{\quad\ \texttt{\# Definir constantes}}

        \State{\(\alpha \gets d/(N\epsilon^{2})\)}
        \For{\(k \in \{1, \dots, K\}\)}
        \State{\(\alpha_{k} \gets D/(\tr(\vPi_{k})\epsilon^{2})\)}
        \State{\(\gamma_{k} \gets \tr(\vPi_{k})/D\)}
        \EndFor
        \Statex{}

        \Statex{\quad\ \texttt{\# Iteración capa por capa de ReduNet}}
        \State{\(\vZ^{1} = \mat{\vz_{1}^1, \dots, \vz_{N}^1} \gets
        \vX\)} \Comment{Inicializar la iteración ReduNet por capas}
        \For{\(\ell \in \{1, \dots, L\}\)}
        \Statex{\qquad\quad\ \texttt{\# Paso 1: Calcular los parámetros de la red}
        \(\vE^{\ell}, \{\vC^{\ell}_{k}\}_{k = 1}^{K}\)}
        \State{\(\vE^{\ell} \gets \alpha\left(\vI + \alpha
        \vZ^{\ell}(\vZ^{\ell})^{\top}\right)^{-1} \in \R^{d \times d}\)}
        \For{\(k \in \{1, \dots, K\}\)}
        \State{\(\vC^{\ell}_{k} \gets \alpha_{k}\left(\vI +
                \alpha_{k}\vZ^{\ell}\vPi_{k}(\vZ^{\ell})^{\top}\right)^{-1}
        \in \R^{d \times d}\)}
        \EndFor
        \Statex{}
        \Statex{\qquad\quad\texttt{\# Paso 2: Actualizar características \(\vZ^{\ell}\)}}
        \For{\(i \in \{1, \dots, N\}\)}
        \Statex{\qquad\qquad\quad\texttt{\# Calcular asignaciones blandas
        \(\hat{\vpi}(\vz^{\ell}_i)\)}}
        \State{\(\hat{\vpi}(\vz^{\ell}_i) \gets \displaystyle
                \softmax(-\lambda[\norm{\vC^{\ell}_{1}\vz^{\ell}_i}_{2},
                \dots, \norm{\vC^{\ell}_{K}\vz^{\ell}_i}_{2}]) \in [0,
        1]^{K}\)}
        \Statex{}
        \Statex{\qquad\qquad\quad\texttt{\# Actualizar
        features \(\vz^{\ell + 1}_i\) from \(\vz^{\ell}_i\)}}
        \State{\(\displaystyle\vz^{\ell + 1}_i \gets \cP_{\bS^{d -
                1}}\rp{\vz^{\ell}_i + \eta \bp{\vE^{\ell}\vz^{\ell}_i -
                        \sum_{k = 1}^{K}\gamma_{k}\hat{\pi}_{k}(\vz^{\ell}_i)
        \vC^{\ell}_{k}\vz^{\ell}_i}} \in \R^{d}\)}
        \EndFor
        \EndFor
        \Statex{}
        \Statex{\quad\ \texttt{\# Devolver todos los parámetros de la red.}}
        \State{\Return{\(\{\vE^{\ell}\}_{\ell = 1}^{L},
                    \{\{\vC^{\ell}_{k}\}_{k = 1}^{K}\}_{\ell = 1}^{L},
        \{\gamma_{k}\}_{k = 1}^{K}\)}}
        \EndProcedure
    \end{algorithmic}
\end{algorithm}

\paragraph{Red profunda para optimizar la reducción de tasas.} Obsérvese que
el incremento está diseñado para emular el ascenso del gradiente para la
reducción de la tasa $\Delta R_\epsilon$. Por lo tanto, al transformar las
características de forma iterativa mediante el proceso anterior, esperamos que la
reducción de la tasa aumente, como veremos en la sección experimental.
Este proceso iterativo, una vez que converge, por ejemplo tras $L$ iteraciones,
proporciona el mapa de características deseado $f(\x, \bm \theta)$ sobre la entrada $\x =
\z^0$, precisamente en forma de una {\em red profunda}, en la que cada
capa tiene la estructura que se muestra en la \Cref{fig:arch-ReduNet} a la izquierda:
\begin{align}\label{eqn:ReduNet}
    f(\x, \bm \theta)\; =&  \;\;f^L \circ f^{L-1} \circ  \cdots
    \circ f^1 \circ
    f^0(\z^0),  \\
    f^\ell(\z^\ell, \bm \theta^\ell) \; \doteq & \;\; \z^{\ell+1}
    = \mathcal{P}_{\mathbb{S}^{n-1}}[\z^{\ell} + \eta\cdot
    g(\z^{\ell}, \bm \theta^{\ell})], \\
    g(\z^{\ell}, \bm \theta^{\ell}) \; =&\;\; \bm E^{\ell}
    \z^{\ell} -  \bm \sigma\big([\bm{C}^{\ell}_{1} \z^{\ell},
    \dots, \bm{C}^{\ell}_{K} \z^{\ell}]\big).
\end{align}
Dado que esta red profunda se deriva de la maximización de la tasa
\textbf{de reducción}, la denominamos \textbf{ReduNet}. Al comparar
la arquitectura de ReduNet con la de redes populares diseñadas empíricamente,
\textbf{ResNet} y \textbf{ResNeXt}, tal como se muestra en la
\Cref{fig:arch-ReduNet}, la similitud resulta casi inquietante.
Desde un punto de vista conceptual, ReduNet también podría utilizarse para justificar la popular
arquitectura de mezcla de expertos (\textbf{MoE}) \cite{MoE}, ya que cada
canal paralelo, $\bm{C}^{\ell}_k$, puede considerarse un «experto»
entrenado para cada clase de objetos.

\begin{figure}[t]
    \centering
    \includegraphics[width=0.45\linewidth]{\toplevelprefix/chapters/chapter5/figs/MoE.png}
    \hspace{5mm}
    \includegraphics[width=0.45\linewidth]{\toplevelprefix/chapters/chapter5/figs/switched-transformer.png}
    \caption{Izquierda: una red profunda con mezcla de expertos (MoE)
        \cite{MoE}. Derecha: un Switch Transformer que promueve la dispersión
    \cite{Fedus-2022}, utilizado para implementar MoE con 1,7 billones de parámetros.}
    \label{fig:enter-label}
\end{figure}

Resumimos el entrenamiento y la evaluación de ReduNet en el
\Cref{alg:redunet_training} y en el \Cref{alg:redunet_evaluation},
respectivamente.  Cabe destacar que todos los parámetros de la red se
construyen explícitamente capa por capa mediante una {\em propagación hacia adelante}.
¡Esta construcción no requiere retropropagación! Las características así 
aprendidas pueden utilizarse directamente para la clasificación, por ejemplo, mediante
un clasificador de subespacio más cercano.

\begin{algorithm}[!htbp]
    \caption{Algoritmo de evaluación para ReduNet}
    \label{alg:redunet_evaluation}
    \begin{algorithmic}[1]
        \Require{Entrada \(\vx \in \R^{D}\), parámetros de la red
            \(\{\vE^{\ell}\}_{\ell = 1}^{L}, \{\{\vC^{\ell}_{k}\}_{k =
            1}^{K}\}_{\ell = 1}^{L}, \{\gamma_{k}\}_{k = 1}^{K}\),
        tasa de aprendizaje \(\lambda\)}
        \Ensure{característica \(\vz^{L + 1}\)}

        \Procedure{ReduNetEvaluation}{$\vx$}
        \Statex{\quad\ \texttt{\# Inicializar
        la iteración por capas de ReduNet}}
        \State{\(\vz^{1} \gets \vx \in \R^{D}\)}
        \Statex{}
        \For{\(\ell \in \{1, \dots, L\}\)}
        \Statex{\qquad\quad\texttt{\# Calcular las asignaciones blandas
        \(\hat{\vpi}(\vz^{\ell})\)}}
        \State{\(\hat{\vpi}(\vz^{\ell}) \gets
                \softmax(-\lambda\mat{\norm{\vC^{\ell}_{1}\vz^{\ell}}_{2},
                \dots, \norm{\vC^{\ell}_{K}\vz^{\ell}}_{2}}) \in [0,
        1]^{K}\)}
        \Statex{}
        \Statex{\qquad\quad\texttt{\# Actualizar
        feature \(\vz^{\ell + 1}\) using \(\vz^{\ell}\)}}
        \State{\(\vz^{\ell + 1} \gets \cP_{\bS^{d - 1}}\rp{\vz^{\ell}
                    + \eta \bp{\vE^{\ell}\vz^{\ell} - \sum_{k =
                        1}^{K}\gamma_{k}\hat{\pi}_{k}(\vz^{\ell})
        \vC^{\ell}_{k}\vz^{\ell}}} \in \R^{d}\)}
        \EndFor
        \Statex{}
        \Statex{\quad\ \texttt{\# Devolver las características de salida}}
        \State{\Return{\(\vz^{L + 1}\)}}
        \EndProcedure
    \end{algorithmic}
\end{algorithm}

\begin{example}
    Para dar una idea de cómo ReduNet transforma las características,
    presentamos un ejemplo sencillo con gaussianas 3D mezcladas y visualizamos
    cómo se transforman las características en la \Cref{fig:redu-3d-gaussian-diagram}.
    Consideremos una mezcla de tres distribuciones gaussianas en $\R^{3}$
    que se proyecta sobre $\mathbb{S}^2$. Primero generamos puntos de datos
    para 3 clases: para $k=1,2,3$, $\bm{X}_{k}=[\bm{x}_{k,1},
    \ldots, \bm{x}_{k,m}] \in \R^{3\times m}$, $\bm{x}_{k,i} \sim
    \mathcal{N}(\bm{\mu}_{k}, \sigma_{k}^{2} \I)$ y ${\pi}(\x_{k,i}) = k$.
    Fijamos $m=500, \sigma_{1}=\sigma_{2}=\sigma_{3}=0.1$ y
    $\bm{\mu}_{1}, \bm{\mu}_{2}, \bm{\mu}_{3} \in \mathbb{S}^2$.
    A continuación, proyectamos todos los puntos de datos en $\mathbb{S}^{2}$, es decir,
    $\bm{x}_{k,i}/\|\bm{x}_{k,i}\|_{2}$.
    Para construir la red (calculando $\bm{E}^{\ell},
    \bm{C}^{\ell}_{k}$ para la capa $\ell$-ésima), establecemos el número de
    iteraciones/capas $L=2,000$, el tamaño del paso $\eta=0.5$ y la precisión
    $\epsilon=0.1$.
    Hacemos esto únicamente para demostrar que nuestro marco permite obtener 
    redes profundas estables incluso con miles de capas.
    En la práctica, puede que no sean necesarias miles de capas y se puede
    detener el proceso cuando la incorporación de nuevas capas empiece a dar resultados cada vez menores.
    Para este ejemplo, bastan unos doscientos niveles.
    Por lo tanto, el objetivo claro de optimización proporciona un criterio natural
    para determinar la profundidad necesaria de la red.

    Como se muestra en la \Cref{fig:redu-3d-gaussian-diagram}, podemos observar
    que, tras la aplicación del mapeo $f(\cdot, \theta)$, las muestras de la
    misma clase se comprimen en gran medida y convergen en un único grupo o clúster,
    y que el ángulo entre dos grupos diferentes es aproximadamente
    $\pi/2$, lo cual concuerda perfectamente con la solución óptima
    $\Z^{\star}$ de la pérdida MCR$^2$ en $\mathbb{S}^2$.
    La pérdida MCR$^2$ de las características en las diferentes capas se puede observar en
    la \Cref{fig:redu-3d-gaussian-diagram}(\textbf{c}). Empíricamente, hemos
    observado que la ReduNet construida es capaz de maximizar la pérdida MCR$^2$
    y converge de manera estable y que las muestras de la misma clase
    convergen en un solo grupo y los puntos de diferentes grupos son
    ortogonales entre sí.\footnote{Debido a la baja dimensión,
        cada grupo se reduce a un subespacio de dimensión \(1\) en la
        esfera (por ejemplo, un solo punto y su antípoda); este es el
        caso extremo de mantener la dimensión dentro del grupo al mínimo. En
    dimensiones superiores, por supuesto, los subespacios estarán menos comprimidos.}
    Además, al tomar nuevas muestras de los mismos
    distribuciones, observamos que las nuevas muestras de la misma clase
    convergen sistemáticamente hacia el mismo centro de clúster que las muestras de entrenamiento.
    \begin{figure}[t]
        \begin{subfigure}[t]{0.32\textwidth}
            \centering
            \includegraphics[width=\textwidth]{\toplevelprefix/chapters/chapter5/figs/scatter3d-X_train.pdf}\vspace{-0.1in}
            \caption{$\bm{X}_{\text{entrenamiento}}$}
        \end{subfigure}
        \hfill
        \begin{subfigure}[t]{0.32\textwidth}
            \centering
            \includegraphics[width=\textwidth]{\toplevelprefix/chapters/chapter5/figs/scatter3d-Z_train.pdf}\vspace{-0.1in}
            \caption{$\bm{Z}_{\text{entrenamiento}}$}
        \end{subfigure}
        \hfill
        \begin{subfigure}[t]{0.32\textwidth}
            \centering
            \includegraphics[width=\textwidth]{\toplevelprefix/chapters/chapter5/figs/scatter3d-loss-traintest.pdf}\vspace{-0.1in}
            \caption{Pérdida (entrenamiento/evaluación)}
        \end{subfigure}
        \vspace{-0.1in}
        \caption{\small Muestras originales y representaciones aprendidas
            para la mezcla de gaussianas en 3D. Visualizamos los puntos de datos
            $\bm{X}$ (antes de la aplicación de $f(\cdot, \theta)$) en
            (\textbf{a}) y las características aprendidas $\bm{Z}$ (después de la aplicación de
            $f(\cdot, \theta)$) en (\textbf{b}) mediante un diagrama de dispersión. En
            cada diagrama de dispersión, cada color representa una clase de
            muestras. En (\textbf{c}), también mostramos los diagramas que representan la
        evolución de los valores de las funciones objetivo.}
        \label{fig:redu-3d-gaussian-diagram}
    \end{figure}
\end{example}

\subsection{Redes convolucionales a partir de la reducción de la tasa en 
invariantes}\label{sec:shift-invariant}

En la sección anterior, derivamos la arquitectura por capas de una
red profunda, la ReduNet, utilizando la optimización desenrollada para el objetivo
de reducción de la tasa.
Concretamente, el término de compresión $R^c_\epsilon(\Z \mid\bm \Pi)$ en
\eqref{eq:mcr-formulation} está diseñado para comprimir representaciones
de la misma clase. Sin embargo, esta formulación no tiene en cuenta
posibles transformaciones o deformaciones del dominio de los datos de entrada. Por
ejemplo, desplazar ligeramente un objeto hacia la derecha no modifica
la etiqueta semántica de una imagen.  En esta sección, demostraremos
cómo se pueden derivar capas convolucionales maximizando un objetivo de
reducción de la tasa que sea invariante ante ciertas deformaciones del dominio,
como las rotaciones y traslaciones de imágenes.

En muchas tareas de agrupamiento o clasificación (como la detección de objetos
en imágenes), consideramos que dos muestras son {\em equivalentes} si
solo difieren en ciertas clases de deformaciones o aumentos del dominio
$\cT = \{\tau \}$. Por lo tanto, solo nos interesan las estructuras de baja dimensión
que sean {\em invariantes} ante tales deformaciones~(es decir, $\x
\in \mathcal{M}$ si y solo si $\tau(\x) \in \mathcal{M}$ para todo $\tau \in \cT$\,),
que se sabe que tienen estructuras geométricas y topológicas sofisticadas
y que pueden ser difíciles de aprender con precisión en la práctica, incluso
con CNNs diseñadas rigurosamente \cite{Cohen-ICML-2016}.
En este marco, esto se puede formular de una manera muy natural: todas las
instancias equivariantes deben integrarse en el mismo subespacio, de modo
que el propio subespacio sea invariante respecto a las transformaciones que
se están considerando.

En muchas aplicaciones, como los datos en serie o los datos de imágenes, el
significado semántico (etiquetas) de los datos es  {\em invariante} ante ciertas
transformaciones $\mathfrak{g} \in \mathbb{G}$ (para algún grupo
$\mathbb{G}$) \cite{CohenW16,deep-sets-NIPS2017}. Por ejemplo, el
significado de una señal de audio es invariante ante un desplazamiento en el tiempo; y la
identidad de un objeto en una imagen es invariante ante una traslación en el
plano de la imagen. Por lo tanto, preferimos que el mapeo de características $f(\x,\bm \theta)$
sea rigurosamente invariante ante tales transformaciones:
\begin{equation}
    \mbox{\em Invarianza de grupo:}\;   f(\x\circ \mathfrak{g}, \bm
    \theta) \sim f(\x,\bm \theta),\ \forall \mathfrak{g} \in \mathbb{G},
\end{equation}
donde «$\sim$» indica que dos características pertenecen a la misma
clase equivalente. Aunque para garantizar la invariancia o la equivariancia,
los operadores convolucionales han sido una práctica común en las redes profundas
\cite{CohenW16}, sigue siendo un reto en la práctica entrenar una
red convolucional (diseñada empíricamente) desde cero que pueda {\em
garantizar} la invariancia incluso ante transformaciones simples como
la traslación y la rotación \cite{azulay2018deep,engstrom2017rotation}.
Un enfoque alternativo consiste en diseñar cuidadosamente los filtros de convolución de
cada capa para garantizar la invariancia traslacional en una amplia gama
de señales, por ejemplo, utilizando ondículas como en ScatteringNet
\cite{scattering-net} y trabajos posteriores \cite{Wiatowski-2018}.
Sin embargo, para garantizar la invariancia frente a señales genéricas, el número
de convoluciones necesarias suele crecer exponencialmente con la profundidad
de la red. Esa es la razón por la que este tipo de red no puede
construirse con tanta profundidad, sino que suele tener solo unas pocas capas.

Ahora mostraremos que el principio MCR$^2$ es compatible con la invariancia
de una manera natural y precisa: solo necesitamos asignar todas las versiones transformadas
$\{\x\circ \mathfrak{g} \mid \mathfrak{g} \in \mathbb G\}$
a la misma clase que los datos $\x$ y mapear todas sus características $\z$
al mismo subespacio $\mathcal S$. Por lo tanto, toda la información equivariante
respecto al grupo se codifica únicamente dentro del subespacio, y cualquier clasificador
definido sobre el conjunto resultante de subespacios será automáticamente
invariante respecto a dichas transformaciones de grupo. Véase la
\Cref{fig:ortho-invariance-diagram} para una ilustración de los
ejemplos de rotación en 1D y traslación en 2D. A continuación, demostraremos rigurosamente
que, cuando el grupo $\mathbb G$ realiza desplazamientos circulares en una dimensión, la
red profunda resultante se convierte naturalmente en una {\em red convolucional
multicanal}.
Dado que la red así construida solo necesita garantizar la invariancia
para los datos dados $\X$ o sus características $\Z$, el número de
convoluciones necesarias se mantiene, de hecho, constante a lo largo de una red muy profunda,
a diferencia de lo que ocurre con ScatteringNet.

\begin{figure}[t]
    \begin{subfigure}[t]{0.4\textwidth}
        \centering
        \includegraphics[width=\textwidth]{\toplevelprefix/chapters/chapter5/figs/ortho_diagram_1d.pdf}
    \end{subfigure}
    \hfill
    \begin{subfigure}[t]{0.4\textwidth}
        \centering
        \includegraphics[width=\textwidth]{\toplevelprefix/chapters/chapter5/figs/ortho_diagram_2d.pdf}
    \end{subfigure}
    \caption{Ilustración de la representación buscada que es
        equivariante/invariante ante la rotación de la imagen (izquierda) o la traslación
        (derecha): todas las imágenes transformadas de cada clase se mapean en el
        mismo subespacio, que es incoherente con respecto a otros subespacios. Las características
        incrustadas en cada subespacio son equivariantes a las transformaciones del grupo,
        mientras que cada subespacio es invariante a dichas
    transformaciones.}
    \label{fig:ortho-invariance-diagram}
\end{figure}

\paragraph{Datos en serie unidimensionales e invariancia de desplazamiento} Para clasificar
datos unidimensionales $\bm x = [x(0), x(1), \ldots, x(D-1)] \in \Re^D$
invariantes bajo traslación, tomamos $\mathbb{G}$ como el grupo de todas las
traslaciones circulares. Cada observación $\bm x_i$ genera una familia $\{
\x_i \circ \mathfrak{g} \, | \, \mathfrak{g} \in \mathbb G \}$ de
copias desplazadas, que son las columnas de la matriz circulante
$\circm(\bm x_i) \in \Re^{D \times D}$ dada por
\begin{equation}
    \circm(\x) \,\doteq\, \left[
        \begin{array}{ccccc} x(0) & x(D-1) & \dots & x(2) & x(1)
            \\ x(1) & x(0) & x(D-1) & \cdots & x(2) \\ \vdots & x(1)
            & x(0) &\ddots & \vdots \\ x(D-2) &  \vdots & \ddots &
            \ddots & x(D-1) \\ x(D-1) & x(D-2) & \dots & x(1) & x(0)
    \end{array} \right]  \in \Re^{D \times D}.
\end{equation}
Remitimos al lector a \cite{Kra2012OnCM} para conocer las propiedades de las matrices circulantes.
Para simplificar, sea $\bm Z^1 \doteq [ \z_{1}^1, \dots,
\z_{N}^1 ] = \X \in \Re^{d \times N}$.\footnote{Una vez más, para simplificar
    el análisis, supongamos por ahora que las características iniciales $\Z^1$ son
    el propio $\X$ y, por lo tanto, tienen la misma dimensión $d$, es decir, $D=d$. Pero
    eso no tiene por qué ser así, ya que pronto veremos que necesitamos
elevar $\X$ a una dimensión superior.} Entonces, ¿qué ocurre si construimos
la ReduNet a partir de sus familias circulantes $\circm(\bm Z^1) = \left[
\circm(\z_{1}^1), \dots, \circm(\z_{N}^1) \right] \in \Re^{d \times
(dN)}$? Es decir, queremos comprimir y proyectar todas estas en el mismo
subespacio mediante la ReduNet.

Obsérvese que ahora la matriz de covarianza de los datos:
\begin{eqnarray}
    \circm(\bm Z^1) \circm(\bm Z^1)^\top
    &=& \left[ \circm(\z_{1}^1), \dots, \circm(\z_{N}^1) \right]
    \left[ \circm(\z_{1}^1), \dots, \circm (\z_{N}^1) \right]^\top \nonumber \\
    &=& \sum_{i =1}^N \circm(\z_{i}^1) \circm(\z_{i}^1)^\top \;\in
    \Re^{d\times d}
\end{eqnarray}
asociada a esta familia de muestras es, {\em de manera automática}, una
matriz circulante (simétrica). Además, dado que la propiedad circulante
se conserva bajo sumas, inversas y productos, las
matrices $\bm E^1$ y $\bm C^1_k$ son también automáticamente matrices circulantes,
cuya aplicación a un vector de características $\bm z \in \Re^d$ puede
implementarse utilizando la convolución circular ``$\circledast$''.
Específicamente, tenemos la siguiente proposición.

\begin{proposition}[Estructuras convolucionales de $\bm E^1$ and $\bm C^1_k$]
    La matriz
    \begin{equation}
        \bm E^1 = \alpha\big(\bm I + \alpha \circm(\bm Z^1)
        \circm(\bm Z^1)^\top \big)^{-1}
    \end{equation}
    es una matriz circulante y representa una convolución circular:
    $$\bm E^1 \z = \bm e_1 \circledast \z,$$
    donde $\bm e_1 \in \Re^d$ es el primer vector columna de $\bm E^1$
    y «$\circledast$» es la convolución circular, definida como
    \begin{equation*}
        (\bm e_1 \circledast \bm z)_{i} \doteq \sum_{j=0}^{d-1}
        e_1(j) x(i+ d-j \,\, \textsf{mod} \,\,d).
    \end{equation*}
    Del mismo modo, las matrices $\bm C^1_k$ asociadas a cualquier subconjunto
    de $\bm Z^1$ también son convoluciones circulares.
    \label{prop:circular-conv-1}
\end{proposition}

No solo los parámetros de la primera capa $\bm E^1$ y $\bm{C}^1_k$ de
la ReduNet se convierten en convoluciones circulantes, sino que también las características de la siguiente capa
siguen siendo matrices circulantes.
Es decir, la transformación incremental de características en
\eqref{eqn:layer-approximate} aplicada a todas las versiones desplazadas de un
$\z^1 \in \Re^d$, dada por
\begin{equation}
    \circm(\z^1) + \eta \cdot \bm E^1 \circm(\z^1) - \eta \cdot \bm
    \sigma \Big([\bm{C}_1^1 \circm(\z^1), \ldots, \bm{C}^1_K
    \circm(\z^1)] \Big),
\end{equation}
es una matriz circulante. Esto implica que no es necesario
construir familias circulantes a partir de las características de la segunda capa, como hicimos
para la primera capa.
Denotando
\begin{equation}
    \z^{2} \propto \z^{1} +\eta \cdot g(\z^{1}, \bm \theta^{1}) =
    \z^1 + \eta \cdot \bm e_{1} \circledast \z^{1} -  \eta \cdot \bm
    \sigma\Big([\bm{c}_{1}^{1} \circledast \z^{1}, \dots,
    \bm{c}^{1}_{K} \circledast \z^{1}]\Big),
    \label{eqn:approximate-convolution}
\end{equation}
las características del siguiente nivel se pueden escribir como
$$\circm(\bm Z^2) = \big[ \circm( \bm z_{1}^1 + \eta g( \bm z_{1}^1,
    \bm \theta^1)), \dots, \circm( \bm z_{N}^1 + \eta g(\bm z_{N}^1, \bm
\theta^1)) \big].$$
Siguiendo con el razonamiento inductivo, vemos que todas las matrices $\bm E^\ell$ y
$\bm C^\ell_k$ basadas en dicho $\circm(\bm Z^\ell)$ son circulantes, y
lo mismo ocurre con todas las características.
Gracias a las propiedades de los datos, ReduNet ha adoptado la forma
de una red convolucional, {\em ¡sin necesidad de elegir explícitamente
esta estructura!}

\paragraph{Una disyuntiva fundamental entre invariancia y dispersión.}
Sin embargo, hay un problema: en general, el conjunto de todas las
permutaciones circulares de un vector $\z$ da lugar a una matriz de rango completo. Es decir, las
$d$ características «aumentadas» asociadas a cada muestra (y, por lo tanto, a cada
clase) suelen abarcar ya todo el espacio $\Re^d$. Por ejemplo,
todas las versiones desplazadas de una función delta $\delta(d)$ pueden generar cualquier
otra señal como su superposición ponderada (densa). El objetivo MCR$^2$
\eqref{eqn:maximal-rate-reduction} no podrá
distinguir las clases como subespacios diferentes.

Un remedio natural consiste en mejorar la separabilidad de los datos «elevando»
la señal original a un espacio de mayor dimensión, por ejemplo,
tomando sus respuestas a múltiples filtros $\bm k_1, \ldots, \bm
k_C \in \Re^d$:
\begin{equation}
    \bm z[c] = \bm k_c \circledast \bm x  =  \circm(\bm k_c) \bm x
    \in \Re^d, \quad c = 1, \ldots, C.
    \label{eqn:lift-1d}
\end{equation}
Los filtros pueden ser filtros prediseñados que promueven la invariancia
,\footnote{Para señales unidimensionales, como el audio, se puede recurrir a la
    transformada de Fourier de corto plazo (STFT) convencional; para imágenes bidimensionales, se
    puede recurrir a las ondas bidimensionales, como en ScatteringNet
\cite{scattering-net}.} o aprendidos de forma adaptativa a partir de los
datos,\footnote{En el caso de los filtros aprendidos, estos pueden aprenderse como
    componentes principales de las muestras, tal como en PCANet
    \cite{chan2015pcanet}, o a partir del aprendizaje de diccionarios de convolución
\cite{li2019multichannel,qu2019nonconvex}.} o seleccionados al azar, tal como
lo hacemos en nuestros experimentos. Esta operación eleva cada señal original
$\x \in \Re^d$ a una característica de $C$ canales, denotada como $\bar{\z}  \doteq
[\z[1], \ldots, \z[C]]^\top \in \Re^{C\times d}$.
Entonces, podemos construir la ReduNet sobre representaciones vectoriales de
$\bar{\z}$, denotadas como
$\vec(\bar\z) \doteq [\z[1]^\top, \ldots, \z[C]^\top] \in \Re^{dC}$.
Entonces, podemos construir la ReduNet sobre representaciones vectoriales de
$\bar{\z}$, denotadas como
$\vec(\bar\z) \doteq [\z[1]^\top, \ldots, \z[C]^\top] \in \Re^{dC}$.
La versión circulante asociada $ \circm(\bar{\z})$ y su 
matriz de covarianza de datos, denotada como $\bar{\bm \Sigma}(\bar\z)$, para todas sus
versiones desplazadas se dan como:
\begin{align}
    \circm(\bar{\z}) &\doteq \left[
        \begin{matrix}
            \circm(\z[1])  \\ \vdots \\ \circm(\z[C])
    \end{matrix} \right] \in \Re^{dC\times d},
    \\
    \bar{\bm \Sigma}(\bar\z) &\doteq
    \left[
        \begin{matrix}
            \circm(\z[1]) \\ \vdots \\ \circm(\z[C])
    \end{matrix} \right]\left[
        \begin{matrix}\circm(\z[1])^\top,\ldots, \circm(\z[C])^\top
    \end{matrix} \right] \in \Re^{dC\times dC},
    \label{eqn:W-multichannel}
\end{align}
donde $\circm(\z[c]) \in \Re^{d\times d}$ con $c \in [C]$ es la
versión circulante del canal $c$-ésimo de la característica $\bar \z$.
Entonces, las columnas de $\circm(\bar\z)$  solo generarán, como máximo, un
subespacio propio de dimensión $d$ en $\Re^{dC}$. Sin embargo, esta simple
operación de elevación (si es lineal) no basta aún para que las clases
sean separables: las características asociadas a otras clases abarcarán el
{\em mismo} subespacio de dimensión $d$. Esto refleja un conflicto fundamental
entre la invariancia y el modelado lineal (de subespacios): {\em no
    se puede esperar que señales desplazadas y superpuestas arbitrariamente pertenezcan
a la misma clase.}

\begin{figure}[t]
    \centerline{
    \includegraphics[width=0.9\textwidth]{\toplevelprefix/chapters/chapter5/figs/sparse-representation.pdf}}
    \caption{Cada señal de entrada $\bm x$ (en este caso, una imagen) puede
        representarse como una superposición de convoluciones dispersas con
    múltiples núcleos $\bm d_c$ en un diccionario $\bm D$.}
    \label{fig:multi-channel-sparse-representation}
\end{figure}
Una manera de resolver este conflicto consiste en aprovechar una
estructura adicional dentro de cada clase, en forma de {\em dispersión}: las señales
dentro de cada clase no se generan como una superposición lineal
arbitraria de ciertos átomos (o motivos) básicos, sino únicamente como combinaciones {\em dispersas}
de estos y sus versiones desplazadas, tal como se muestra en la
\Cref{fig:multi-channel-sparse-representation}. Más concretamente, sea
$\bm D_k = [\bm d_{k,1}, \ldots, \bm d_{k,c}]$ una matriz con un
conjunto de átomos asociados a la clase $k$, también conocido como
diccionario; entonces, cada señal $\x$ de esta clase se genera en forma dispersas como:
\begin{equation}
    \x = \bm d_{k,1} \circledast z_1 + \ldots + \bm d_{k,c}
    \circledast z_c = \circm(\bm{D}_k)\z,
\end{equation}
para algún vector disperso $\z$. Las señales de las diferentes clases se
generan entonces mediante diferentes diccionarios cuyos átomos (o motivos) son
incoherentes entre sí. Debido a esta incoherencia, es poco probable que las señales de una clase
sean representadas de forma dispersa por átomos de cualquier otra clase.
Por lo tanto, todas las señales pueden representarse como
\begin{equation}
    \x = \big[\circm(\bm{D}_1), \circm(\bm{D}_2), \ldots,
    \circm(\bm{D}_K)\big] \bar \z,
\end{equation}
donde $\bar \z$ es un vector disperso.\footnote{Cabe señalar que desde hace tiempo se han propuesto y utilizado modelos similares de representación dispersa
    con fines de clasificación en aplicaciones como el reconocimiento facial,
    demostrando una eficacia excelente
    \cite{Wright:2009,wagner2012toward}. Recientemente, el modelo de codificación escasa convolucional
    ha sido propuesto por \cite{papyan2017convolutional} como un
    marco para interpretar las estructuras de las redes convolucionales 
profundas.} Existe una amplia bibliografía sobre cómo aprender los
diccionarios de dispersión más compactos y óptimos a partir de datos de muestra, por ejemplo,
\cite{li2019multichannel,qu2019nonconvex}, y posteriormente resolver el
problema inverso y calcular el código disperso asociado $\z$ o $\bar
\z$. Estudios recientes de \cite{qu2020nonconvex,Qu2020Geometric} incluso
demuestran que, en condiciones generales, el problema del aprendizaje de diccionarios de convolución
puede resolverse de manera eficaz y eficiente.

No obstante, para tareas como la clasificación, no nos
interesa necesariamente el diccionario óptimo exacto ni el
código disperso exacto para cada señal individual. Lo que nos
interesa principalmente es que, en conjunto, el conjunto de códigos dispersos de cada clase sea
suficientemente separable del de las demás clases. Bajo el
supuesto del modelo generativo disperso, si los núcleos de convolución
$\{\bm k_c\}_{c=1}^C$  se ajustan bien a la «transpuesta» o
«inversa» de los diccionarios de dispersión anteriores $\bm D = [\bm D_1,
\ldots, \bm D_K]$, también conocidos como los {\em filtros de análisis}
\cite{Cosparse-Nam,Analysis-Filter}, las señales de una clase solo
tendrán respuestas altas a un pequeño subconjunto de esos filtros y respuestas bajas
a los demás (debido a la suposición de incoherencia).
No obstante, en la práctica, a menudo basta con un número suficientemente grande de, digamos,
$C$ filtros aleatorios $\{\bm k_c\}_{c=1}^C$ para garantizar que las
características extraídas de los $C$ canales
\begin{equation}
    \big[\bm k_1 \circledast \x, \bm k_2 \circledast \x, \ldots, \bm
    k_C \circledast \x\big]^\top = \big[\circm(\bm k_1) \x, \ldots,
    \circm(\bm k_C) \x \big]^\top \in \Re^{C\times d}
\end{equation}
para las diferentes clases presentan distintos patrones de respuesta ante distintos
filtros, lo que permite separarlas \cite{chan2015pcanet}.

Por lo tanto, en nuestro marco, en gran medida el número de canales
(o la amplitud de la red) desempeña realmente el papel de  {\em
recurso estadístico}, mientras que el número de capas (la profundidad de la
red) desempeña el papel de {\em recurso computacional}. La
teoría de la detección compresiva caracteriza con precisión cuántas
mediciones se necesitan para preservar las estructuras intrínsecas
de baja dimensión (incluida la separabilidad) de los datos
\cite{Wright-Ma-2022}. %

\begin{figure}[t]
    \centerline{
    \includegraphics[width=0.95\textwidth]{\toplevelprefix/chapters/chapter5/figs/sparse-lifting.pdf}}
    \caption{Estimación del código disperso $\bar{\bm z}$ de una señal de entrada
        $\bm x$ (en este caso, una imagen) mediante convoluciones con
    varios núcleos $\bm k_c$ y, a continuación, la dispersión.}
    \label{fig:multi-channel-sparse-lifting}
\end{figure}
Las respuestas multicanal $\bar \z$ deben ser dispersas. Por lo tanto, para
aproximar el código disperso $\bar \z$, podemos aplicar un umbral no lineal 
elemento por elemento que promueva la dispersión, por ejemplo, $\bm \tau(\cdot)$, a
las salidas del filtro anterior, estableciendo que las respuestas bajas (por ejemplo, 
con valor absoluto inferior a $\epsilon$) o negativas sean cero:
\begin{equation}
    \bar \z \doteq \bm \tau \left( \big[\circm(\bm k_1) \x, \ldots,
    \circm(\bm k_C) \x \big]^\top \right) \in \Re^{C \times d}.
    \label{eqn:sparse-lifting}
\end{equation}
La \Cref{fig:multi-channel-sparse-lifting} ilustra las ideas básicas.
Se puede consultar \cite{Analysis-Filter} para un estudio más sistemático
sobre el diseño del operador de umbralización de dispersión. Sin embargo,
aquí no nos interesa tanto obtener los mejores códigos dispersos,
siempre y cuando los códigos sean suficientemente separables. Por lo tanto, el operador no lineal
$\bm \tau$ puede elegirse simplemente como un umbral blando o una función ReLU.
Se puede suponer que estas características presumiblemente dispersas $\bar \z$ se encuentran en una
subvariedad (no lineal) de dimensión inferior de $\mathbb{R}^{dC}$, la cual
puede linealizarse y separarse de las otras clases mediante capas
ReduNet posteriores, como se ilustra más adelante en la \Cref{fig:learn-to-classify-diagram}.

La red ReduNet construida a partir de la versión circulante de estas características multicanal
$\bar\Z \doteq [\bar \z_1, \ldots, \bar \z_N] \in \Re^{C
\times d \times N}$, es decir, $\circm(\bar\Z) \doteq [ \circm(\bar\z_1),
\dots, \circm(\bar\z_N)] \in \Re^{dC \times dN}$, conserva las buenas
propiedades de invariancia descritas anteriormente: los operadores lineales, ahora
denominados $\bar{\bm E}$ y $\bar{\bm C}_k$, siguen siendo circulantes por bloques,
y representan convoluciones circulares unidimensionales multicanal.
Concretamente, obtenemos el siguiente resultado.
\begin{proposition}[Estructuras de convolución multicanal de $\bar{\bm
    E}$ y $\bar{\bm C}_k$]
    La matriz
    \begin{equation}
        \label{eq:def-E-bar}
        \bar{\bm E} \doteq \alpha\left(\bm I + \alpha\,
        \circm(\bar\Z) \circm(\bar\Z)^\top \right) ^{-1}
    \end{equation}
    es circulante por bloques, es decir,
    \begin{equation}
        \bar{\bm E} =
        \left[
            \begin{matrix}
                \bar{\bm E}_{1, 1} & \cdots & \bar{\bm E}_{1, C}\\
                \vdots & \ddots & \vdots \\
                \bar{\bm E}_{C, 1} & \cdots & \bar{\bm E}_{C, C}\\
        \end{matrix}\right] \in \Re^{dC \times dC},
    \end{equation}
    donde cada $\bar{\bm E}_{c, c'}\in \Re^{d \times d}$ es una
    matriz circulante. Además, $\bar{\bm E}$ representa una
    convolución circular multicanal, es decir, para cualquier señal multicanal
    $\bar\z \in \Re^{C \times n}$ tenemos
    \begin{equation}
        \bar{\bm E} \cdot \textsf{vec}(\bar\z) = \textsf{vec}( \bar{\bm
        e} \circledast \bar\z).
    \end{equation}
    En lo anterior, $\bar{\bm e} \in \Re^{C \times C \times d}$ es un
    núcleo convolucional multicanal donde $\bar{\bm e}[c, c'] \in
    \Re^{d}$ es el primer vector columna de $\bar{\bm E}_{c, c'}$
    y $\bar{\bm e} \circledast \bar\z \in \Re^{C \times d}$ es la
    convolución circular multicanal definida como
    \begin{equation}
        (\bar{\bm e} \circledast \bar\z)[c] \doteq \sum_{c'=1}^C
        \bar{\bm e}[c, c'] \circledast \bar{\z}[c'], \quad \forall c
        = 1, \ldots, C.
    \end{equation}
    Del mismo modo, las matrices $\bar{\bm C}_k$ asociadas a cualquier
    subconjunto de $\bar{\bm Z}$ son también convoluciones circulares multicanal.
    \label{prop:multichannel-circular-conv-1}
\end{proposition}
Según la \Cref{prop:multichannel-circular-conv-1}, la ReduNet invariante ante traslaciones
es, por su propia construcción, una red convolucional profunda para señales unidimensionales multicanal. Obsérvese que, incluso si los kernels de elevación iniciales están
separados \eqref{eqn:sparse-lifting}, la inversa de la matriz  en
\eqref{eq:def-E-bar} para calcular $\bar{\bm E}$ (de manera similar para
$\bar{\bm C_k}$) introduce «interferencia» entre todos los $C$ canales. %
Por lo tanto, estas convoluciones multicanal, en general, {\em no}
son separables en profundidad, a diferencia de las redes Xception~\cite{Xception} que
en su momento se propusieron para simplificar las redes neuronales convolucionales
multicanal.\footnote{Queda por determinar qué estructuras adicionales en los
datos darían lugar a convoluciones separables en profundidad.}

\begin{remark}[Reducción de la complejidad computacional en el dominio de frecuencias]
    El cálculo de $\bar{\bm E}$ en \eqref{eq:def-E-bar} requiere
    invertir una matriz de tamaño $dC \times dC$, lo que en general tiene
    una complejidad de $O(d^3C^3)$.  Sin embargo, al aprovechar el hecho de que una
    matriz circulante puede diagonalizarse mediante la matriz de la Transformada de Fourier Discreta
    (DFT), la complejidad puede reducirse significativamente.
    Como se muestra en \cite{chan2021redunet}, para calcular $\bar{\bm
    E}$ y $\bar{\bm C}_k \in \Re^{dC \times dC}$, solo necesitamos
    calcular en el dominio de la frecuencia la inversa de los bloques $C\times C$
    dentro de $d$ veces, por lo que la complejidad total es de $O(dC^3)$.

\end{remark}

\paragraph{Arquitectura general de la red y comparación.}
A partir de la derivación anterior, observamos que, para encontrar una
representación discriminativa lineal (LDR) para múltiples clases de
señales o imágenes que sea invariante ante traslaciones, la codificación
dispersa, una arquitectura multicapa con convoluciones multicanal, diferentes
funciones de activación no lineales y el cálculo del espectro se convierten en componentes {\em
necesarios} para alcanzar el objetivo de manera eficaz y
eficiente. La \Cref{fig:learn-to-classify-diagram} ilustra el
proceso general de aprendizaje de dicha representación mediante la reducción invariante de tasas 
sobre los códigos dispersos de entrada.

\begin{figure}[t]
    \centering
    \includegraphics[width=0.98\linewidth]{\toplevelprefix/chapters/chapter5/figs/learn_to_classify_diagram_updated.png}
    \caption{El proceso general para clasificar señales multiclase
        con invariancia al desplazamiento: elevación multicanal,  codificación dispersa,
        seguida de una ReduNet de convolución multicanal para la reducción de la tasa
        invariante. Estos componentes son {\em necesarios} para
        mapear señales multiclase invariantes al desplazamiento a subespacios incoherentes (lineales)
        como una LDR. Cabe señalar que las arquitecturas de la mayoría de las
        redes neuronales profundas modernas se asemejan a este proceso. El LDR así aprendido
    facilita tareas posteriores, como la clasificación.}
    \label{fig:learn-to-classify-diagram}
\end{figure}

\begin{example}[Clasificación invariante de dígitos]
    A continuación, presentamos los resultados empíricos de ReduNet en el
    aprendizaje de características invariantes a la \textit{rotación} en el conjunto de datos real
    MNIST de 10 clases.
    Imponemos una cuadrícula polar sobre la imagen $\bm{x}\in\mathbb{R}^{H\times
    W}$, cuyo centro geométrico coincide con el centro de la cuadrícula polar 2D
    (tal como se ilustra en la \Cref{fig:samples-invariant-1d-mnist-diagram}).
    Para cada radio $r_i$, $i \in [C]$, podemos muestrear $\Gamma$ píxeles
    con respecto a cada ángulo $\gamma_l =l\cdot({2\pi}/\Gamma)$ con
    $l \in [\Gamma]$.
    Entonces, dada una imagen de muestra $\bm{x}$ del conjunto de datos, representamos
    la imagen en coordenadas polares (muestreadas) como una señal multicanal
    $\bm{x}_p \in\R^{\Gamma\times C}$.
    El objetivo aquí es aprender una representación invariante ante rotaciones,
    es decir, esperamos aprender $f(\cdot, \theta)$ tal que
    $\{f(\bm{x}_p \circ \mathfrak{g}, \theta)\}_{\mathfrak{g}
    \in\mathbb{G}}$ se encuentren en el mismo subespacio, donde $\mathfrak{g}$ es
    el desplazamiento cíclico en el ángulo polar.
    Utilizamos $N=100$ muestras de entrenamiento ($10$ de cada clase) y establecemos
    $\Gamma=200$, $C=15$ para el muestreo polar.
    Al realizar el muestreo anterior en coordenadas polares, podemos
    obtener la matriz de datos $\bm{X}_p \in \mathbb{R}^{(\Gamma\cdot C) \times N}$.
    Para ReduNet, establecemos el número de capas/iteraciones $L=40$,
    la precisión $\epsilon=0.1$ y el tamaño del paso $\eta=0.5$. Antes de la primera
    capa, realizamos una elevación de la entrada mediante una convolución
    circulante unidimensional con 20 núcleos gaussianos aleatorios de tamaño 5.

    \begin{figure}[t]
        \begin{subfigure}[t]{0.3\textwidth}
            \centering
            \includegraphics[width=\textwidth]{\toplevelprefix/chapters/chapter5/figs/1d-rotation.pdf}
            \caption{$\bm{X}_{\text{rotación}}$}
        \end{subfigure}
        \hfill
        \begin{subfigure}[t]{0.3\textwidth}
            \centering
            \includegraphics[width=\textwidth]{\toplevelprefix/chapters/chapter5/figs/mnist1d_img0.pdf}
            \caption{$\bm{Z}_{\text{rotación}}$}
        \end{subfigure}
        \hfill
        \begin{subfigure}[t]{0.3\textwidth}
            \centering
            \includegraphics[width=\textwidth]{\toplevelprefix/chapters/chapter5/figs/mnist1d_img1.pdf}
            \caption{Pérdida}
        \end{subfigure}
        \caption{\small Ejemplos de imágenes giradas de los dígitos MNIST,
            cada una 18$^{\circ}$. (\textbf{Izquierda}) Diagrama de la
            representación en coordenadas polares;  (\textbf{Derecha}) Imágenes giradas
        de los dígitos «0» y «1».}
        \label{fig:samples-invariant-1d-mnist-diagram}
    \end{figure}

    Para evaluar la representación aprendida, cada muestra de entrenamiento se
    aumenta con 20 versiones rotadas de sí misma, cada una desplazada con
    un paso de 10. Calculamos las similitudes coseno entre las entradas de entrenamiento aumentadas de $m \times
    20$ entradas de entrenamiento aumentadas $\bm{X}_{\text{rotation}}$ y los
    resultados se muestran en la \Cref{fig:redu-invariant-1d-mnist-diagram}
    (\textbf{a}).
    Comparamos las similitudes coseno entre las características aprendidas de
    todas las versiones aumentadas, es decir,
    $\bar{\bm{Z}}_{\text{rotation}}$ y resumimos los resultados en la
    \Cref{fig:redu-invariant-1d-mnist-diagram} (\textbf{b}).
    Como vemos, la ReduNet invariante a la rotación así construida es capaz
    de mapear los datos de entrenamiento (así como todas sus versiones rotadas)
    de las 10 clases diferentes en 10 subespacios casi ortogonales.
    Es decir, los subespacios aprendidos son verdaderamente invariantes ante
    la transformación de desplazamiento en ángulo polar. A continuación, extraemos al azar
    otras $100$ muestras de prueba y aplicamos el mismo procedimiento de aumento.
    En la \Cref{fig:redu-invariant-1d-mnist-diagram} (\textbf{c}),
    visualizamos la pérdida MCR$^{2}$ en la representación de la capa $\ell$
    de la ReduNet sobre los conjuntos de datos de entrenamiento y de prueba.
    A partir de estos resultados, podemos observar que la ReduNet construida es
    efectivamente capaz de maximizar la pérdida MCR$^{2}$, así como de generalizar
    a los datos de prueba.

    \begin{figure}[t]
        \begin{subfigure}[t]{0.3\textwidth}
            \centering
            \includegraphics[width=\textwidth]{\toplevelprefix/chapters/chapter5/figs/mnist1d-heatmap-X_translate_train_all.pdf}
            \caption{$\bm{X}_{\text{rotación}}$}
        \end{subfigure}
        \hfill
        \begin{subfigure}[t]{0.3\textwidth}
            \centering
            \includegraphics[width=\textwidth]{\toplevelprefix/chapters/chapter5/figs/mnist1d-heatmap-Z_translate_train_all.pdf}
            \caption{$\bm{Z}_{\text{rotación}}$}
        \end{subfigure}
        \hfill
        \begin{subfigure}[t]{0.32\textwidth}
            \centering
            \includegraphics[width=\textwidth]{\toplevelprefix/chapters/chapter5/figs/mnist1d-loss-traintest.pdf}
            \caption{Pérdida}
        \end{subfigure}
        \caption{\small (a) y (b) son mapas de calor de la similitud coseno
            entre los datos de entrenamiento rotados $\bm{X}_{\text{rotation}}$ y
            las características aprendidas $\bar{\bm{Z}}_{\text{rotation}}$ para
            la invariancia a la rotación. (d) muestra las pérdidas MCR$^2$
        de entrenamiento y validación por capas.}
        \label{fig:redu-invariant-1d-mnist-diagram}
    \end{figure}

\end{example}

\section{Transformadores de white-box mediante optimización 
desenrollada}\label{sec:chap4-white-box-transformer}
Como hemos visto en la sección anterior, utilizamos el problema de la
clasificación para ofrecer una interpretación rigurosa de las principales
características arquitectónicas de redes profundas populares como
ResNet y las CNN: cada capa de estas redes puede considerarse como
una imitación de un paso de gradiente que aumenta el objetivo de
reducción de la tasa (o ganancia de información). Esta perspectiva también lleva a un
hecho algo sorprendente: los parámetros y operadores  de las
capas de una red tan profunda, la ReduNet, pueden  calcularse de forma
puramente hacia adelante.

A pesar de la importancia teórica y conceptual de ReduNet,
hay varios factores que limitan su aplicabilidad práctica. En primer lugar, como hemos
comentado anteriormente, el costo computacional de calcular los
operadores matriciales en cada capa de forma directa puede ser muy
elevado. En segundo lugar, es posible que los operadores así calculados no sean tan eficaces a la hora de
optimizar el objetivo y que se necesiten miles de iteraciones 
(y, por lo tanto, de capas). Como hemos visto en la \Cref{sec:LISTA} para LISTA, estos
dos problemas pueden resolverse permitiendo optimizar esos operadores
y haciendo que sean aprendibles mediante retropropagación.\footnote{O, tal vez,
mediante una combinación de optimización hacia adelante y hacia atrás.}

El contexto de clasificación supervisada en el que se desarrolló ReduNet
también presenta ciertas limitaciones. En la práctica, es posible que una imagen no
pertenezca a una sola clase, ya que puede contener varios objetos. Por lo tanto,
sería más general suponer que las diferentes regiones de la imagen
pertenecen a  diferentes modelos de baja dimensión (por ejemplo, una distribución gaussiana o un
subespacio). Como veremos, tal generalización daría lugar a una arquitectura
tanto sencilla como general que unifica las operaciones de reducción de tasa y
de eliminación de ruido que hemos visto en el capítulo anterior.
Además, la arquitectura así obtenida se asemeja a la popular
arquitectura de Transformador.

\subsection{Optimización desenrollada para la reducción dispersa de la tasa}

En los últimos años, a medida que la cantidad y la calidad de los datos
disponibles para entrenar redes profundas han aumentado de manera espectacular,
el campo ha pasado de considerar las entradas como «átomos» a considerarlas como
«secuencias de tokens», en las que cada token forma parte de la entrada global.
Este cambio de perspectiva ofrece ventajas tanto filosóficas como prácticas:
desde el punto de vista filosófico, nos permite medir y modelar cada
parte de la entrada y sus interacciones; desde el punto de vista práctico, nos permite
utilizar redes profundas similares, \textit{modelos secuencia-a-secuencia}
como los omnipresentes Transformadores, para aprender todo tipo de
distribuciones de datos. Ahora nos corresponde comprender el aprendizaje de representaciones
en el modelo «centrado en tokens» de datos como secuencias de tokens,
y hacer frente a los retos mencionados anteriormente (como, por ejemplo, en el
contexto del marco de reducción de tasa, la falta de etiquetas para cada token).

Formalmente, sea \(\X = \mat{\x_{1}, \dots, \x_{N}} \in \bR^{D \times
N}\) un conjunto de variables aleatorias que representan nuestra fuente de datos. En tareas de visión,
cada \(\x_{i} \in \bR^{D}\) se interpreta como un
\textit{token}, que suele corresponder a un fragmento de imagen. En
tareas de lenguaje, cada \(\x_{i} \in \bR^{D}\) se interpreta como una
\textit{incrustación de token}, es decir, una representación vectorial continua de
un token discreto, como una palabra o una subpalabra.\footnote{Con un ligero
    desvío terminológico, a lo largo de este capítulo nos referiremos tanto a los tokens discretos como a sus
    incrustaciones asociadas simplemente como «tokens» por
razones de conveniencia.}  %
Los \(\x_{i}\) pueden tener estructuras de correlación arbitrarias. Utilizamos
\(\Z = \mat{\z_{1}, \dots, \z_{N}} \in \bR^{d \times N}\) para denotar
las variables aleatorias que definen nuestras representaciones, donde \(\z_{i}
\in \bR^{d}\) es la representación del token correspondiente \(\x_i
\in \bR^{D}\).%

\begin{remark}
    En los transformadores, cada muestra de entrada se convierte normalmente en una
    secuencia de {\em tokens}. Un token es una unidad básica de información
    derivada de la entrada sin procesar: en el procesamiento del lenguaje natural,
    los tokens suelen ser palabras o subpalabras; en la visión artificial,
    corresponden a fragmentos de imagen; y en otras modalidades, pueden
    representar intervalos de tiempo, ubicaciones espaciales u otras unidades
    específicas del dominio. Una {\em incrustación de token} es una representación vectorial continua
    de un token que sirve como entrada para un
    transformador. Asigna a cada token un punto en un espacio de alta dimensión,
    lo que permite al modelo procesar entradas simbólicas mediante
    cálculos numéricos.
    Una {\em representación de token} es un vector que codifica la
    información semántica o estructural de un token, y suele ser generada
    por las capas intermedias o finales de un transformador. Estas
    representaciones están diseñadas para capturar características significativas de
    la entrada que resultan útiles para tareas posteriores, tales como
    la clasificación, la generación o la regresión. Consulte la
    \Cref{sec:contrastive_learning} para obtener más detalles sobre este
    vocabulario y cómo se relaciona con las implementaciones.
\end{remark}

\paragraph{Objetivo para el aprendizaje de una representación estructurada y compacta.}
Inspirándonos en el análisis previo sobre la reducción de la tasa
(\Cref{sec:chap4-white-box-model-via-unrolling}), proponemos que el
objetivo del aprendizaje de representaciones
en este modelo centrado en tokens para nuestros datos sea encontrar una
correspondencia de características \(f
\colon \X \in \bR^{D \times N} \to \Z\in \bR^{d \times N}\) que
transforme los tokens de entrada \(\{\bm x_i\}_{i=1}^N \subset \R^D\) con una
distribución potencialmente no lineal y multimodal a representaciones de tokens (por tramos)
\textit{linealizadas y compactas} \(\{\bm
z_i\}_{i=1}^N \subset \bR^{d}\). Si bien la distribución {\em conjunta}
de las representaciones de los tokens \(\{\z_{i}\}_{i = 1}^{N}\) puede ser
sofisticada (y específica para cada tarea), sostenemos además que es
razonable y práctico
exigir que la distribución {\em marginal} objetivo de las representaciones
individuales de los tokens sea altamente comprimida y estructurada,
lo que la hace apta para una codificación compacta. En particular, requerimos que la
distribución sea \textit{una mezcla de distribuciones gaussianas de baja dimensión (digamos \(K\))
}, de tal manera que la \(k\)-ésima gaussiana tenga media \(\Zero \in
\bR^{d}\), covarianza \(\vSigma_{k} \succeq \Zero \in \bR^{d \times
d}\), y soporte generado por la base ortonormal \(\vU_{k} \in
\bR^{d \times p}\).
Denotamos por $\vU_{[K]} = \{\vU_k\}_{k=1}^K$ el conjunto de bases de
todas las distribuciones gaussianas.

Como es habitual, para aprender buenas representaciones, queremos maximizar
la \textit{ganancia de información} de las representaciones finales de los tokens, o
más concretamente, maximizar una noción adecuada de la reducción de la tasa
(véase \Cref{subsec:MCR2}), es decir,
\begin{align}\label{eq:rate reduction}
    \mathrm{max}_{\bm Z \in \R^{d\times N}}\ \Delta R_{\epsilon}(\Z
    \mid \vU_{[K]}) \doteq R_{\epsilon}(\Z) - R^c_{\epsilon}(\Z \mid \vU_{[K]}).
\end{align}
Aquí, el primer término $R_{\epsilon}$ es una estimación de la tasa de codificación con pérdida
para el conjunto completo de representaciones de tokens. Más
concretamente, si consideramos las representaciones de tokens $\{\bm
z_i\}_{i=1}^N$ como muestras independientes e idénticamente distribuidas de una única distribución gaussiana de media cero,
su tasa de codificación con pérdida, sujeta a una precisión de cuantificación $\epsilon
> 0$, viene dada por
\begin{equation}\label{eq:coding_rate}
    R_{\epsilon}(\Z) \doteq \frac{1}{2}\textrm{logdet}\left(\I +
    \frac{d}{N\epsilon^{2}}\Z^{\top}\Z\right) =
    \frac{1}{2}\textrm{logdet}\left(\I +
    \frac{d}{N\epsilon^{2}}\Z\Z^{\top}\right).
\end{equation}
El segundo término $R_{\epsilon}^c$ es una estimación de la tasa de codificación con pérdida
bajo el libro de códigos $\bm U_{[K]}$, que se expresa como
\begin{equation}
    \scalebox{0.95}{\(\displaystyle R_{\epsilon}^c(\Z \mid \vU_{[K]}) \doteq
            \sum_{k=1}^{K}R_{\epsilon}(\vU_k^{\top} \Z) =
            \frac{1}{2}\sum_{k=1}^{K}\log\det\left(\I +
    \frac{p}{N\epsilon^{2}}(\vU_k^{\top}\Z)^{\top}(\vU_k^{\top}\Z)\right).\)}
    \label{eq:def-mcr-Rc}
\end{equation}

\begin{remark}
    La expresión \eqref{eq:def-mcr-Rc} para la tasa de codificación puede
    considerarse como una generalización de la tasa de codificación \(R_\epsilon^{c}\)
    utilizada en el objetivo de reducción de tasa original \eqref{eq:MCRc}. En
    particular, el objetivo original se define con respecto a un
    conjunto de etiquetas de pertenencia conocidas \(\{\vPi_{k}\}\) específicas de la
    realización de datos particular \(\vX\). Por el contrario, el objetivo actual
    se define con respecto a los subespacios \(\vU_{[K]}\),
    que son independientes de cualquier realización concreta, pero de los que
    se supone que admiten la distribución de las representaciones de tokens.
    Supongamos que una representación de token $\bm z_i$ pertenece a un
    subespacio $\bm U_k$ y que estos subespacios son aproximadamente
    ortogonales entre sí, es decir, $\bm U_k^\top \bm U_l \approx \bm
    0$ para todo $k \neq l$. Entonces, se puede verificar que las proyecciones
    $\bm U_k\bm U_k^\top \bm z_i  = \bm z_i$ y  $\bm U_l\bm
    U_l^\top \bm z_i \approx \bm 0$ para todo $l \neq k$. Estas
    proyecciones ortogonales sirven efectivamente como etiquetas de pertenencia implícitas,
    identificando el subespacio al que pertenece cada representación de token.
\end{remark}

\begin{remark}
    Cabe destacar que este conjunto de requisitos para el aprendizaje de representaciones
    concuerda bien con dos hipótesis empíricas bien establecidas sobre
    la estructura de las representaciones de tokens en los modelos entrenados que
    utilizan la omnipresente arquitectura de redes neuronales de Transformadores: la
    «hipótesis de la representación lineal»
    \citep{jiang2024origins,park2023linear} y la «hipótesis de la superposición»
    » \citep{elhage2022toy,yun2021transformer}. La hipótesis de la
    representación lineal postula que las representaciones de tokens en
    los modelos de Transformador se encuentran en subespacios lineales de baja dimensión que
    codifican características semánticas. De manera similar, la hipótesis de la superposición
    sugiere que estas representaciones pueden expresarse, de manera aproximada,
    como una combinación lineal dispersa de estos vectores de características. Como veremos en
    la \Cref{fig:crate_segmentation_specialization} más adelante en este
    capítulo, la base del subespacio $\bm U_k$ puede interpretarse como un conjunto
    de características semánticas, donde cada característica corresponde a un
    aspecto específico del significado del token. Las representaciones de los tokens se
    expresan entonces, aproximadamente, como combinaciones lineales dispersas de
    estas bases del subespacio, capturando los componentes semánticos esenciales
    del token e ignorando las dimensiones irrelevantes.
\end{remark}

\begin{figure}[t!]
    \centering
    \includegraphics[width=0.95\textwidth]{\toplevelprefix/chapters/chapter5/figs/coding-transform.png}
    \caption{ \small\textbf{Comparación de tres conjuntos de
        representaciones mediante la reducción de la tasa y la dispersión. Cada $S_i$
        representa un subespacio lineal y el número de bolas azules
        representa la diferencia entre las tasas de codificación $\Delta
        R_{\epsilon}(\vZ \mid \vU_{[K]}) = R_\epsilon(\vZ) -
        R^c_\epsilon(\vZ \mid \vU_{[K]})$.}
    }
    \label{fig:sparse-rate-reduction-diagram}
\end{figure}

\paragraph{Reducción dispersa de tasas.} Tenga en cuenta que el objetivo de reducción de la tasa
\eqref{eq:rate reduction} es invariante ante rotaciones conjuntas arbitrarias
de las representaciones y los subespacios. En particular,
la optimización del objetivo de reducción de la tasa puede no conducir de forma natural a
representaciones alineadas con los ejes (es decir, \textit{dispersas}). {Por ejemplo,
    considere los tres conjuntos de representaciones aprendidas en la
    \Cref{fig:sparse-rate-reduction-diagram}.} Por lo tanto, nos gustaría
transformar las representaciones (y sus subespacios
de soporte) de tal manera que las representaciones $\vZ$ acaben siendo
dispersas\footnote{Concretamente, con pocas entradas distintas de cero.} con respecto
a las coordenadas estándar del espacio de representaciones resultante {como
en la \Cref{fig:sparse-rate-reduction-diagram}(c)}. 
Desde el punto de vista computacional, podemos combinar los dos objetivos anteriores en un único
objetivo de optimización:
\begin{equation}
    \max_{f \in \mathcal{F}}\ [\Delta R_{\epsilon}(\Z \mid \vU_{[K]})
    - \lambda \|\bm{Z}\|_0] \qquad \text{s.t.}\ \Z = f(\X),
    \label{eq:sparse-rr}
\end{equation}
donde $\mathcal{F}$ denota una clase general de funciones y la norma $\ell_0$
$\|\Z\|_0$ promueve la dispersión de las representaciones finales de los tokens
\(\Z = f(\X)\).%

En la práctica, la norma $\ell_0$ suele sustituirse por la norma $\ell_1$
para mejorar la manejabilidad computacional y permitir el uso de técnicas
de optimización convexa \cite{Wright-Ma-2022}. Motivados por esto, simplificamos el Problema
\eqref{eq:sparse-rr} en consecuencia, lo que da lugar a una formulación que
se mantiene fiel al objetivo de dispersión original, al tiempo que resulta más
apta para algoritmos eficientes, tal como se muestra a continuación:
\begin{equation}
    \begin{aligned}
        \max_{f \in \mathcal{F}}\ [\Delta R_{\epsilon}(\Z \mid
        \vU_{[K]}) - \lambda \|\bm{Z}\|_1]  \qquad \text{s.t.}\ \Z = f(\X),
        \label{eq:sparse-rr-1}
    \end{aligned}
\end{equation}
Aunque se trate de un ligero uso impropio de la terminología, a menudo nos referimos a esta función objetivo
también como la \textit{reducción dispersa de tasa }.

\paragraph{Arquitectura de red de white-box mediante optimización desenrollada.}

Aunque resulte fácil de formular, la optimización de cada término del objetivo anterior
plantea un reto computacional \cite{Wright-Ma-2022}. Por lo tanto,
es natural adoptar un enfoque de aproximación que realice la
transformación global $f$ para optimizar 
\eqref{eq:sparse-rr} mediante una concatenación de múltiples, digamos $L$,
operaciones simples \textit{incrementales y locales} $f^\ell$ que empujan
la distribución de la representación hacia la distribución del modelo parsimonioso
deseada:
\begin{equation}
    f\colon \X \xrightarrow{\hspace{1mm} f^{\mathrm{pre}} \hspace{1mm}}
    \Z^1 \rightarrow \cdots \rightarrow \Z^\ell
    \xrightarrow{\hspace{1mm} f^{\ell} \hspace{1mm}} \Z^{\ell+1}
    \rightarrow  \cdots \xrightarrow{\hspace{1mm} f^{L}} \Z^{L+1} = \Z,
    \label{eq:incremental}
\end{equation}
donde $f^{\mathrm{pre}}: \bR^{D} \rightarrow \bR^{d}$ es la aplicación de preprocesamiento
que transforma cada token de entrada $\x_{i} \in \bR^{D}$ en las
representaciones iniciales de los tokens $\z_{i}^{1} \in \bR^{d}$.
Cada \textit{mapeo hacia adelante} incremental $\Z^{\ell + 1} =
f^\ell(\Z^\ell)$, o una «capa», transforma la distribución de tokens
para \textit{optimizar} el objetivo de reducción dispersa de tasa anterior
\eqref{eq:sparse-rr}, condicionado a la distribución de su entrada $\Z^\ell$.

\begin{remark}
    A diferencia de otros enfoques de optimización desenrollada, como el
    ReduNet (véase la \Cref{sec:chap4-white-box-model-via-unrolling}), nosotros
    \textit{modelamos explícitamente} la distribución de $\Z^\ell$ en cada
    capa, por ejemplo, como una mezcla de subespacios lineales o generada de forma dispersa
    a partir de un diccionario. Los parámetros del modelo se aprenden a partir de los datos
    (por ejemplo, mediante la \textit{retropropagación} con entrenamiento de extremo a extremo).
    Esta separación entre la «optimización» hacia adelante y el
    «aprendizaje» hacia atrás aclara la función matemática de cada capa como un operador
    que transforma la distribución de su entrada, mientras que la distribución
    de la entrada, a su vez, es modelada (y posteriormente aprendida) por los
    parámetros de la capa.
\end{remark}

A continuación, mostramos cómo derivar estas operaciones incrementales y locales
desde una perspectiva de optimización desenrollada para resolver el problema
\eqref{eq:sparse-rr-1}. Una vez que decidimos utilizar un enfoque incremental
para optimizar el problema
\eqref{eq:sparse-rr-1}, existen diversas opciones posibles para
lograr la optimización. Dado un modelo para $\Z^\ell$, por ejemplo, una mezcla
de subespacios $\vU_{[K]}$, optamos por un método de \textit{minimización
alternada} en dos pasos con una sólida base conceptual. Primero,
\textit{comprimimos} los tokens $\vZ^{\ell}$ mediante un descenso de gradiente para
minimizar el término de tasa de codificación $R^c_\epsilon (\vZ \mid
\vU_{[K]}^\ell)$. Específicamente, tomamos un paso de gradiente en
$R^c_\epsilon$ con una tasa de aprendizaje $\kappa$ de la siguiente manera:
\begin{align}\label{eq:Z l+1/2}
    \bm Z^{\ell+1/2} = \bm Z^\ell - \kappa \nabla_{\bm Z}
    R^c_\epsilon (\vZ \mid \vU_{[K]}^\ell).
\end{align}
A continuación, \textit{raleamos} los tokens comprimidos, generando
\(\vZ^{\ell + 1}\) mediante un paso de gradiente proximal adecuadamente relajado para
minimizar el término restante $\lambda \norm{\vZ}_{1} -
R_{\epsilon}(\vZ)$. Como argumentaremos en detalle más adelante, podemos encontrar
dicho $\bm Z^{\ell+1}$ resolviendo un problema de presentación dispersa con
respecto a un diccionario ortogonal completo $\bm D^\ell$:
\begin{equation}\label{eq:Z l+1}
    \vZ^{\ell+1} = \argmin_{{\vZ}}  \bigg\{\lambda \norm{\vZ}_1 +
    \frac{1}{2}\norm{\vZ^{\ell + 1/2} - \vD^{\ell} {\vZ}}_F^2\bigg\}.
\end{equation}
A continuación, ofrecemos detalles técnicos sobre cada uno de los dos
pasos anteriores y derivamos actualizaciones eficientes para su implementación.

\paragraph{La autoatención como descenso por gradiente en la tasa de codificación de las representaciones 
de tokens.} En el primer paso \eqref{eq:Z l+1/2}, el gradiente
de la tasa de codificación \(\nabla_{\bm Z} R^c_\epsilon\) resulta costoso de
calcular, ya que implica \(K\) inversas de matrices distintas, una para cada
uno de los \(K\) subespacios con base \(\vU_{k}^{\ell}\):
\begin{equation}
    \nabla_{\bm Z} R_{\epsilon}^c(\vZ \mid \vU_{[K]})
    = \frac{p}{N\epsilon^2}\sum_{k=1}^K \vU_k\vU_k^\top\vZ\Big(\I +
    \frac{p}{N\epsilon^2}(\vU_k^\top\vZ)^\top(\vU_k^\top\vZ)\Big)^{-1}.
    \label{eq:rate-gradient}
\end{equation}
Ahora demostramos que este gradiente puede aproximarse de forma natural
utilizando el denominado operador de autoatención de subespacios de múltiples cabezas (MSSA),
que tiene una forma funcional similar a la del operador de autoatención de múltiples cabezas
\citep{vaswani2017attention} con \(K\) cabezas (es decir, una para
cada subespacio, procedente de cada inversa de matriz). Aquí, calculamos por aproximación
el gradiente \eqref{eq:rate-gradient} utilizando la serie de Neumann de primer orden
(véase \Cref{ex:neumannn}):\footnote{Esta aproximación solo
    tiene sentido cuando \(\vZ\) tiene una norma pequeña, lo cual, en la práctica, puede
    lograrse mediante estrategias de normalización (como la denominada «normalización
    por capas»). No la incluimos aquí por razones de
    conveniencia notacional, aunque el operador de normalización por capas correspondiente sí
    se incluye en la arquitectura final.}
\begin{align}\label{eq:gd_rc_neumann}
    &\nabla_{\bm Z} R_{\epsilon}^{c}(\vZ \mid \vU_{[K]}) \\
    &\approx \frac{p}{N\epsilon^2} \sum_{k =
    1}^{K}\vU_{k}\vU_{k}^\top\vZ\left(\vI - \frac{p}{N\epsilon^2}
    (\vU_{k}^\top\vZ)^\top(\vU_{k}^\top\vZ)\right) \\
    &= \frac{p}{N\epsilon^2} \left(\sum_{k = 1}^{K}
    \vU_{k}\vU_{k}^\top\right)\vZ -  \left(
    \frac{p}{N\epsilon^2}\right)^2\sum_{k = 1}^{K}
    \vU_{k}(\vU_{k}^\top\vZ)(\vU_{k}^\top\vZ)^\top(\vU_{k}^\top\vZ).
\end{align}
En la aproximación anterior, calculamos la similitud entre las representaciones proyectadas
de los tokens $\{\bm U_k^\top\bm z_i\}_{i = 1}^{n}$ mediante una
autocorrelación entre las características proyectadas, expresada como
$(\vU_{k}^\top\vZ)^\top(\vU_{k}^\top\vZ)$. Aunque se normalicen todas las características de los tokens
dentro de \(\vZ\), la norma de esta matriz puede crecer
arbitrariamente\footnote {Dada la elección práctica de
    \(\epsilon^{2} = p/N\), es decir, se utiliza un código más preciso para
incorporar más información de más muestras.} a medida que aumenta la longitud de la secuencia,
lo que hace que nuestra aproximación de Neumann
sea cada vez más deficiente (ya que el gradiente verdadero no crece arbitrariamente en
tamaño con la longitud de la secuencia). Para resolver este problema práctico y
generalizar nuestra formulación, podemos utilizar en su lugar un núcleo arbitrario para
medir la similitud. Un núcleo especialmente adecuado
que se ajusta a nuestra intención de realizar la autorregresión únicamente
con respecto a los tokens que pertenecen al mismo subespacio es el núcleo local
(Nadaraya-Watson), que implica la estimación de pertenencia mediante
softmax \citep{chaudhari2021attentive}, es decir, el cálculo de la similitud (relativa)
como $\softmax((\vU_{k}^\top\vZ)^\top(\vU_{k}^\top\vZ))$. 
Supongamos ahora que los subespacios se expanden «tanto como sea posible», es
decir, que todas las bases de los subespacios $\bm U_{[K]}$ juntas generan todo el espacio.
Entonces, tenemos que $\sum_{k = 1}^{K} \vU_{k}\vU_{k}^\top = \bm I$. Por lo tanto,
\eqref{eq:gd_rc_neumann} se convierte en
\begin{align}\label{eq:grad Rc}
    \nabla_{\bm Z} R_{\epsilon}^{c}(\vZ \mid \vU_{[K]})
    \approx  \frac{p}{N\epsilon^2} \vZ -  \left(
    \frac{p}{N\epsilon^2}\right)^2 \operatorname{MSSA}(\vZ^{\ell} \mid
    \vU_{[K]}^{\ell}),
\end{align}
donde MSSA se define mediante un operador SSA de la siguiente manera:
\begin{align}
    & \mathrm{SSA}(\vZ \mid \vU_{k})
    \doteq (\vU_{k}^\top
    \vZ)\softmax((\vU_{k}^\top\vZ)^\top(\vU_{k}^\top\vZ)),
    \ \forall k \in [K], \label{eq:SSA} \\
    & \mathrm{MSSA}(\vZ \mid \vU_{[K]})
    \doteq \frac{p}{N\epsilon^2} \mat{\vU_{1}, \dots,
    \vU_{K}}\mat{\mathrm{SSA}({\vZ \mid \vU_{1}}) \\ \vdots
    \\ \mathrm{SSA}({\vZ \mid \vU_{K}})}.\label{eq:Multi-Head-SSA}
\end{align}
Al sustituir \eqref{eq:grad Rc} en \eqref{eq:Z l+1/2}, se obtiene que
se puede aproximar de forma natural mediante
\begin{equation}
    \vZ^{\ell + 1/2} = \left(1 - \frac{\kappa p}{N\epsilon^2}\right)
    \vZ^{\ell} + \frac{\kappa p}{N\epsilon^2}
    \mathrm{MSSA}(\vZ^{\ell} \mid \vU_{[K]}^{\ell}).
    \label{eq:gd-mcr-parts}
\end{equation}

\begin{remark}
    El operador SSA de \eqref{eq:SSA} se asemeja al
    \textit{operador de atención} de un transformador típico
    \citep{vaswani2017attention}, salvo que aquí los
    operadores lineales de valor, clave y consulta  están todos configurados para ser \textit{los
    mismos} que la base del subespacio, es decir, $\vV_{k} = \vK_{k} =
    \bm{Q}_{k} = \vU_k^*$. Por lo tanto, denominamos
    $\mathrm{SSA}({\spcdot\mid\vU_k}): \bR^{d\times n} \rightarrow
    \bR^{p\times n}$  un operador de \textbf{S}ubspace
    \textbf{S}elf-\textbf{A}ttention (SSA). A continuación, el operador
    MSSA completo de la ecuación \eqref{eq:Multi-Head-SSA}, definido formalmente como
    \(\mathrm{MSSA}({\spcdot \mid \vU_{[K]}}) \colon \bR^{d \times n}
    \to \bR^{d \times n}\) y denominado operador del \textbf{M}ulti-Head
    \textbf{S}ubspace \textbf{S}elf-\textbf{A}ttention (MSSA), agrega 
    las salidas de los cabezales de atención promediándolas
    mediante pesos dependientes del modelo, un concepto similar al popular
    operador de autoatención multi-cabezal en las redes transformadoras
    existentes. El paso de gradiente global \eqref{eq:gd-mcr-parts}
    se asemeja a la autoatención de múltiples cabezas implementada con una conexión de salto
    en los transformadores.
\end{remark}

\paragraph{MLP como descenso por gradiente proximal para la codificación dispersa de
representaciones de tokens.} Para el segundo paso de la minimización
alternante, necesitamos minimizar $\lambda \norm{\vZ}_{1} -
R_{\epsilon}(\vZ)$. Obsérvese que el gradiente \(\nabla
R_{\epsilon}(\vZ)\) implica una inversa de matriz y, por lo tanto, el
gradiente proximal ingenuo (véase \Cref{subsec:pgd}) para optimizar este problema
se vuelve inmanejable en problemas a gran escala. %
Por lo tanto, adoptamos un enfoque diferente para equilibrar
la diversidad representativa y la dispersión: postulamos un
diccionario (completo) incoherente u ortogonal $\vD^{\ell} \in \bR^{d
\times d}$ y buscamos ralear las iteraciones intermedias $\vZ^{\ell
+ 1/2}$ con respecto a \(\vD^{\ell}\). Es decir, $\vZ^{\ell + 1/2}
\approx \vD^{\ell} \vZ^{\ell + 1}$ donde $\vZ^{\ell + 1}$ es más
disperso; es decir, es una \textit{codificación esparsa} de \(\vZ^{\ell +
1/2}\). El diccionario \(\vD^{\ell}\) se utiliza para reducir la densidad de todos los tokens
simultáneamente.
Según la suposición de incoherencia, tenemos que $(\vD^{\ell})^\top(\vD^{\ell})
\approx \vI$. Por lo tanto, a partir de \eqref{eq:coding_rate} tenemos que
\begin{equation}
    R_{\epsilon}(\vZ^{\ell + 1/2}) \approx
    R_{\epsilon}(\vD^{\ell}\vZ^{\ell + 1}) \approx R_{\epsilon}(\vZ^{\ell + 1}).
\end{equation}
Para resolver $\lambda \norm{\vZ}_{1} - R_{\epsilon}(\vZ)$, optimizamos
el siguiente problema
\begin{align*}
    \vZ^{\ell + 1} \approx \argmin_{\vZ}  \norm{\vZ}_1 \quad
    \mbox{sujeto a} \quad \vZ^{\ell + 1/2} = \vD^{\ell}\vZ.
\end{align*}
{El programa de representación dispersa anterior suele resolverse
    relajándolo a un programa convexo sin restricciones, conocido como LASSO
\citep{Wright-Ma-2022}:}
\begin{equation}
    \vZ^{\ell + 1} \approx \argmin_{\vZ} \left[\lambda \norm{\vZ}_1 +
    \frac{1}{2}\norm{\vZ^{\ell + 1/2} - \vD^{\ell} \vZ}_F^2 \right].
\end{equation}
En nuestra implementación, también añadimos una restricción no negativa a
$\vZ^{\ell + 1}$ y resolvemos el LASSO no negativo correspondiente:
\begin{equation}
    \vZ^{\ell + 1} \approx \argmin_{\vZ \geq \bm 0}
    \left[\lambda\norm{\vZ}_1 + \frac{1}{2}\norm{\vZ^{\ell + 1/2} -
    \vD^{\ell} \vZ}_{F}^{2}\right].
    \label{eq:sparse-nonnegative}
\end{equation}
A continuación, optimizamos de forma incremental la \Cref{eq:sparse-nonnegative}
realizando un paso desenrollado de {\em descenso de gradiente proximal}, conocido como
paso ISTA \citep{beck2009fast}, para obtener la actualización:
\begin{align}\label{eq:ista-block}
    \vZ^{\ell + 1}
    &= \mathrm{ISTA}({\vZ^{\ell + 1/2} \mid \vD^{\ell}}), \\
    \text{where} \quad \mathrm{ISTA}({\vZ \mid \vD})
    &\doteq \operatorname{ReLU}(\vZ - \eta \vD^\top(\vD\vZ - \vZ) -
    \eta \lambda \bm{1}).
\end{align}

\begin{figure}
    \centering
    \includegraphics[width=0.75\textwidth]{\toplevelprefix/chapters/chapter5/figs/crate_encoder_architecture.pdf}
    \caption{\small \textbf{Una capa de la arquitectura del codificador CRATE.} 
        La arquitectura completa no es más que una concatenación de
        dichas capas, con un tokenizador inicial, un módulo de preprocesamiento
    y un módulo final específico para la tarea (es decir, un módulo de clasificación).}
    \label{fig:crate_backbone}
\end{figure}

\subsection{Arquitectura general del transformador de white-box: CRATE}

Ahora diseñamos una arquitectura de transformador de white-box, denominada Coding
RATE Transformer (CRATE), mediante el desenrollamiento de las actualizaciones anteriores. Al
combinar los dos pasos anteriores \eqref{eq:gd-mcr-parts} y \eqref{eq:ista-block}:
\begin{enumerate}
    \item Compresión local de tokens dentro de una muestra hacia una
        estructura de mezcla de subespacios, lo que da lugar al bloque de autoatención de subespacios
        multicabezal -- \text{MSSA};
    \item Dispersión global de conjuntos de tokens en todas las muestras
        mediante codificación dispersa, lo que da lugar al bloque de dispersión --
        \text{ISTA};
\end{enumerate}
podemos obtener la siguiente capa de transformadores basada en la reducción de la tasa,
que se ilustra en la \Cref{fig:crate_backbone},
\begin{equation}
    \vZ^{\ell+1/2} \doteq \vZ^{\ell} + \text{MSSA}(\vZ^{\ell} \mid
    \vU_{[K]}^{\ell}),
    \qquad
    \vZ^{\ell+1}\doteq \text{ISTA}(\vZ^{\ell+1/2} \mid \bm D^\ell).
\end{equation}
Al componer varias capas de este tipo siguiendo la construcción incremental
de nuestra representación en \eqref{eq:incremental}, obtenemos una
arquitectura de transformador de white-box que transforma los tokens de datos
hacia una unión compacta y dispersa de subespacios incoherentes, donde
$f^{\pre}: \bR^{D \times N} \rightarrow \bR^{d \times N}$ es la
mapeo de preprocesamiento que transforma los tokens de entrada $\vX \in
\bR^{D \times N}$ en representaciones de la primera capa $\vZ^{1} \in \bR^{d
\times N}$. En la 
\Cref{fig:crate-diagram} se muestra un esquema general de esta arquitectura.

\begin{figure}[t!]
    \centering
    \includegraphics[width=\textwidth]{\toplevelprefix/chapters/chapter5/figs/CRATE_fig1_patches.pdf}
    \vspace{-0.1in}
    \caption{
        \textbf{El «bucle principal» del diseño de redes profundas de white-box
        de CRATE.}
        Tras codificar los datos de entrada como una secuencia de tokens $\vZ^0$,
        CRATE construye una red profunda que transforma los
        datos en una configuración canónica de subespacios de baja dimensión
        mediante sucesivas   {\textit{\textbf{compresiones}}}
        frente a un modelo local de la distribución, generando
        $\vZ^{\ell+1/2}$ y      {\textit{\textbf{raleza}}}
        frente a un diccionario global, generando $\vZ^{\ell+1}$.
        Al apilar repetidamente estos bloques y entrenar los parámetros del modelo
        mediante retropropagación, se obtiene una representación potente e
        interpretable de los datos.
    }
    \label{fig:crate-diagram}
\end{figure}

\begin{remark}[\textbf{El papel del paso hacia adelante y de la propagación 
    hacia atrás}]\label{sub:forward_backward}
    A diferencia de otros enfoques de optimización desenrollada, como el
    ReduNet \cite{chan2021redunet}, \textit{modelamos explícitamente} la
    distribución de cada $\vZ^\ell$ y $\vZ^{\ell + 1/2}$ en cada
    capa, ya sea mediante una mezcla de subespacios lineales o generada de forma dispersa
    a partir de un diccionario. Presentamos la interpretación
    de que, en cada capa \(\ell\), las bases aprendidas para los subespacios
    \(\vU_{[K]}^{\ell}\) y los diccionarios aprendidos \(\vD^{\ell}\)
    sirven conjuntamente como un \textit{libro de códigos} o \textit{filtro de análisis}
    que codifica y transforma las representaciones intermedias en
    cada capa \(\ell\). Dado que la distribución de entrada a la capa
    \(\ell\) se modela primero mediante \(\vU_{[K]}^{\ell}\) y luego
    se transforma mediante \(\vD^{\ell}\), la distribución de entrada a cada
    capa es diferente, por lo que necesitamos un libro de códigos independiente en cada
    capa para obtener la codificación más parsimoniosa. Los parámetros de
    estos libros de códigos (es decir, las bases del subespacio y los diccionarios),
    que hasta ahora se suponía que se conocían perfectamente, en realidad se aprenden
    a partir de los datos (por ejemplo, mediante \textit{retropropagación} dentro
    del entrenamiento de extremo a extremo).

    Por lo tanto, nuestra metodología presenta una clara separación conceptual
    entre la «optimización» hacia adelante y el «aprendizaje» hacia atrás para
    la red neuronal profunda de white-box así obtenida. Es decir, en su
    paso hacia adelante, interpretamos cada capa como un operador que,
    condicionado por un modelo aprendido (es decir, un libro de códigos) para la
    distribución de su entrada, transforma esta distribución hacia una
    representación más parsimoniosa. En su propagación hacia atrás,
    el libro de códigos de este modelo, para la distribución de la entrada en
    cada capa, se actualiza para ajustarse mejor a una determinada relación (supervisada)
    entre entrada y salida, tal y como se ilustra en la
    \Cref{fig:forward-backward}. Esta interpretación conceptual
    implica un cierto agnosticismo de las representaciones del modelo
    respecto a la tarea y la función de pérdida concretas; en particular, muchos tipos
    de tareas y funciones de pérdida garantizarán que los modelos de cada capa estén
    bien ajustados, lo que asegura que el modelo produzca representaciones parsimoniosas.
\end{remark}

Hay tres preguntas lógicas que podemos plantearnos sobre CRATE:
\begin{itemize}
    \item Estos modelos CRATE construidos, ¿aprenden características útiles mediante
        la optimización de la reducción dispersa de tasa?
    \item Estos modelos CRATE construidos, ¿aprovechan sus características de alta calidad
        para obtener buenos resultados en tareas de predicción? 
    \item Estos modelos CRATE construidos, ¿muestran fenómenos interesantes,
        por ejemplo, relacionados con la interpretabilidad, como resultado de
        su diseño simplificado? 
\end{itemize}
Sostenemos que la respuesta a cada pregunta es «sí». En primer lugar, en
la \Cref{fig:crate_compression_sparsity} se analiza cómo la red
optimiza cada término de la reducción de la tasa de dispersión; cada capa
disminuye de manera consistente la medida de compresión y las medidas de dispersión
de las características durante la pasada hacia adelante. En segundo lugar, en la
\Cref{tab:crate_comparison_with_sota} mostramos que CRATE presenta
una precisión de clasificación «top-1», comparable, en configuraciones de aprendizaje por transferencia,
a la arquitectura de red neuronal Vision Transformer (ViT)
\citep{dosovitskiy2020image} en un número similar de
parámetros, con un comportamiento de escalabilidad prometedor: obtenemos
mejoras constantes y consistentes al aumentar el tamaño del modelo.
En general, CRATE logra un rendimiento prometedor en conjuntos de datos a gran escala del mundo real
mediante la implementación directa de nuestra arquitectura basada en principios.
Ofrecemos más detalles sobre la implementación y
el análisis de los resultados experimentales en la \Cref{sec:image_classification}.

\begin{figure}[t]
    \begin{subfigure}[t]{0.48\textwidth}
        \centering
        \includegraphics[width=\textwidth]{\toplevelprefix/chapters/chapter5/figs/forward.png}
        \caption{\bf Paso hacia adelante}
    \end{subfigure}
    \hfill
    \begin{subfigure}[t]{0.48\textwidth}
        \centering
        \includegraphics[width=\textwidth]{\toplevelprefix/chapters/chapter5/figs/backward.png}
        \caption{\bf Retropropagación}
    \end{subfigure}
    \caption{\small {\bf Los roles del paso hacia adelante y de la propagación hacia atrás
        en CRATE}. (a) Dados subespacios y
        diccionarios fijos $\{(\bm U_{[K]}^{\ell}, \bm D^{\ell})\}_{\ell=1}^L$,
        cada capa optimiza de manera incremental la reducción de la tasa de dispersión de las
        representaciones en el paso hacia adelante; (b) La retropropagación aprende
        subespacios y diccionarios $\{(\bm U_{[K]}^{\ell}, \bm
    D^{\ell})\}_{\ell=1}^L$ a partir de los datos de entrenamiento. }
    \label{fig:forward-backward}
\end{figure}

\begin{figure}[t!]
    \centering
    \begin{subfigure}{0.49\textwidth}
        \includegraphics[width=\textwidth]{\toplevelprefix/chapters/chapter5/figs/crate_compression_layer.pdf}
        \caption{}
    \end{subfigure}
    \hfill
    \begin{subfigure}{0.49\textwidth}
        \includegraphics[width=\textwidth]{\toplevelprefix/chapters/chapter5/figs/crate_sparsity_layer.pdf}
        \caption{}
    \end{subfigure}
    \caption{\small\textbf{Optimización de las medidas de compresión y dispersión
        mediante la pasada directa de CRATE.} Medimos la
        compresión \(R_{\epsilon}^{c}(\vZ^{\ell + 1/2})\) (a) y la dispersión
        \(\norm{\vZ^{\ell + 1}}_{1}\) (b) en la salida de cada capa. Observamos
        que ambas medidas disminuyen a lo largo del
        paso hacia adelante, tal como lo predice la teoría, y que casi todas las capas
        disminuyen monótonamente cada medida. (La excepción es la dispersión de la última
    capa, ya que las características densas son importantes para la predicción.)}
    \label{fig:crate_compression_sparsity}
\end{figure}

\begin{table*}[t!]
    \centering
    \caption{\small \textbf{Precisión de clasificación «top-1» de CRATE
            en diversos conjuntos de datos con diferentes escalas de modelo
        cuando se entrenan previamente en ImageNet-1K.} Para ImageNet-1K/ImageNet-1K
        ReaL, evaluamos directamente la precisión «top-1». Para otros
        conjuntos de datos, utilizamos modelos entrenados previamente en ImageNet como
        inicialización y evaluamos el rendimiento del aprendizaje por transferencia
    mediante el ajuste fino.}
    \label{tab:crate_comparison_with_sota}
    \small
    \setlength{\tabcolsep}{13.6pt}
    \resizebox{0.98\textwidth}{!}{
        \begin{tabular}{@{}lcccc|cc@{}}
            \toprule
            \textbf{Model} & CRATE{-T}  &
            CRATE{-S} & CRATE{-B} &
            CRATE{-L} & { ViT-T} &  {
            ViT-S } \\
            \midrule
            \midrule
            \# parámetros & 6.09M & 13.12M & 22.80M & 77.64M & {
            5.72M} & { 22.05M} \\
            \midrule
            ImageNet-1K & 66.7 & 69.2 & 70.8 & 71.3 & {
            71.5} & { 72.4} \\
            ImageNet-1K ReaL & 74.0 & 76.0 & 76.5 & 77.4 & {
            78.3 } & { 78.4} \\
            \midrule
            CIFAR10 & 95.5 & 96.0 & 96.8 & 97.2 & {
            96.6} & { 97.2} \\
            CIFAR100 & 78.9 & 81.0 & 82.7 & 83.6 & {
            81.8} & { 83.2}\\
            Oxford Flowers-102 & 84.6 & 87.1 & 88.7 & 88.3 & {
            85.1} & { 88.5}\\
            Oxford-IIIT-Pets & 81.4 & 84.9 & 85.3 & 87.4 & {
            88.5} & { 88.6} \\
            \bottomrule
    \end{tabular}}
\end{table*}

\begin{figure}[t!]
    \centering
    \includegraphics[width=\textwidth]{\toplevelprefix/chapters/chapter5/figs/crate_semantic_heads.pdf}
    \caption{\textbf{Visualización de los mapas de atención de CRATE: ¿qué
        analiza el modelo para realizar una clasificación?} Visualizamos
        los mapas de atención de la última capa, que constituyen la última fuente de
        interacción token a token en la red y, por lo tanto, en cierto
        sentido, la descripción más detallada de qué secciones de la imagen de entrada
        utiliza el modelo para realizar una predicción. Observamos un
        comportamiento \textit{extremadamente interesante}: cuando se entrenan exclusivamente
        en tareas de clasificación, los mapas de atención identifican los objetos
        en primer plano de la imagen; además, ¡se \textit{especializan de forma natural
        en} diferentes tipos de objetos en primer plano! Este
        resultado empírico nos ofrece una interpretación de los subespacios de las capas posteriores
        como si cada uno correspondiera a un conjunto casi disjunto de
    \textit{conceptos}.}
    \label{fig:crate_segmentation_specialization}
\end{figure}

\section{Variantes de arquitecturas profundas por diseño}
\label{sec:chap4-derive-white-box-transformer-variants}

Hasta ahora, esperamos haber aportado pruebas convincentes de que
el papel de las redes profundas (populares) consiste en aplicar ciertos
algoritmos de optimización para minimizar la tasa de codificación (o maximizar la
ganancia de información) de las representaciones aprendidas. Sin embargo, los lectores
familiarizados con los métodos de optimización habrán notado que
las arquitecturas mencionadas (ReduNet o CRATE) corresponden a
técnicas de optimización bastante básicas. Es posible que tengan mucho margen
de mejora en cuanto a eficiencia o eficacia. Además, si
consideramos que el marco teórico propuesto para interpretar las redes profundas
es correcto, no solo debería ayudar a explicar las arquitecturas existentes,
sino que también debería guiarnos en el desarrollo de arquitecturas más eficientes y
eficaces. En esta sección, demostramos que así es:
las nuevas arquitecturas resultantes no solo son totalmente
interpretables, sino que también ofrecen una corrección garantizada y una eficiencia mejorada.

\subsection{Arquitectura de transformadores de solo atención} \label{sub:aot}

En esta subsección, planteamos y respondemos la pregunta: «¿Cuál es la
red neuronal más simple posible que se adapte razonablemente bien y
comprima la representación de forma demostrable?». Según nuestra teoría previa,
la respuesta debería ser una pila de capas MSSA, concebidas como
operadores de compresión con pérdida. Esta arquitectura de red es lo suficientemente
simple como para permitir analizar el rendimiento de compresión o
de eliminación de ruido de forma rigurosa y sistemática. Con este
fin, formalicemos nuestro modelo de subespacios: las representaciones iniciales de los tokens
$\bm Z^{1}$ se obtienen mediante muestreo de una mezcla ruidosa de
distribuciones gaussianas de bajo rango, definidas en subespacios generados por \(\vU_{k} \in
\R^{d \times p}\), de la siguiente manera: dividimos los \(N\) índices de tokens
en subconjuntos conocidos \(C_{1}, \dots, C_{K}\), de tal manera que:
\begin{equation}\label{def:MoG}
    \bm z_i = \underbrace{\bm U_{k} \bm a_i}_{signal} +
    \underbrace{\sum_{j \neq k}^K \bm U_j \bm e_{i,j}}_{
    noise},\ \forall i \in C_k,
\end{equation}
donde $\bm{a}_i \overset{i.i.d.}{\sim}
\mathcal{N}(\bm{0},\bm{I}_{p_k})$ y $\bm{e}_{i,j}
\overset{i.i.d.}{\sim} \mathcal{N}(\bm{0},\delta^2\bm{I}_{p_j})$
para todo $i \in C_k$ y $k \in [K]$, $\{\bm{a}_i\}$ y
$\{\bm{e}_{i,j}\}$ son todas independientes.

\paragraph{Un operador de reducción de ruido para representaciones de tokens.} Ahora,
mostramos que el operador MSSA (véase \eqref{eq:Multi-Head-SSA}) puede
eliminar el ruido de forma incremental de las representaciones de tokens generadas a partir del
modelo anterior. Concretamente, consideramos un operador MSSA generalizado: para cada
$\ell \in [L]$,
\begin{align}\label{eq:MSSA}
    \bm Z^{\ell+1} =  \bm Z^{\ell} + \eta \sum_{k=1}^K \bm U_k\bm
    U_k^T \bm Z^{\ell} \varphi(\bm Z^{\ell^T}\bm U_k\bm
    U_k^T\bm Z^{\ell}),
\end{align}
donde $\eta > 0$ es el tamaño del paso y $\varphi$ es una función
que actúa sobre cada columna de la matriz de entrada, como softmax,
ReLU elemento por elemento, etc. Para simplificar nuestra teoría, suponemos que los
subespacios de la \Cref{def:MoG} son ortogonales entre sí, es decir, $\bm
U_k^T\bm U_j = \bm 0$ para todo $k \neq j$. Cabe señalar que esta suposición
no es restrictiva, ya que en espacios de alta dimensión, los subespacios aleatorios
de baja dimensión son incoherentes entre sí con alta
probabilidad, es decir, $\bm U_k^T\bm U_j \approx \bm 0$
\citep{Wright-Ma-2022}.\footnote{Se pueden generalizar directamente
    nuestros resultados a subespacios no ortogonales, con un análisis ligeramente más
sofisticado.}

\begin{figure}[t]
    \begin{subfigure}[t]{0.45\textwidth}
        \centering
        \includegraphics[width=\textwidth]{\toplevelprefix/chapters/chapter5/figs/SNR1.png}
        \caption{Nivel de ruido $\delta = 0.2$}
    \end{subfigure}
    \hfill
    \begin{subfigure}[t]{0.45\textwidth}
        \centering
        \includegraphics[width=\textwidth]{\toplevelprefix/chapters/chapter5/figs/SNR2.png}
        \caption{Nivel de ruido $\delta = 0.5$}
    \end{subfigure}
    \caption{{\bf Desempeño en la eliminación de ruido del transformador
        basado únicamente en atención.} Aquí, obtenemos muestras de las representaciones iniciales de los tokens a partir de
        una mezcla de gaussianas de bajo rango en la \Cref{def:MoG}. A continuación, aplicamos
        \eqref{eq:MSSA} para actualizar las representaciones de los tokens y presentamos la
    relación señal-ruido (SNR) en cada capa.}
    \label{fig:MSSA}
\end{figure}

Ahora, supongamos que las columnas de $\bm Z_k^{\ell}$ denotan las representaciones de tokens
del subespacio $k$-ésimo en la capa $\ell$-ésima. Para
cuantificar la capacidad de eliminación de ruido, definimos la razón señal-ruido
(SNR) para cada bloque de las representaciones de tokens en la
capa $\ell$-ésima de la siguiente manera:
\begin{align}\label{def:SNR}
    \mathrm{SNR}(\bm Z_k^{\ell}) \doteq  \frac{\|\bm U_k\bm U_k^T\bm
    Z_k^{\ell} \|_F}{\|(\bm I - \bm U_k\bm U_k^T)\bm Z_k^{\ell}
    \|_F},\quad \forall k \in [K].
\end{align}
Para simplificar nuestro análisis, suponemos que $\abs{C_{1}} = \cdots =
\abs{C_{K}} = N/K$ y
\begin{align}\label{eq:orth}
    \begin{bmatrix}
        \bm U_1 & \dots & \bm U_K
    \end{bmatrix} \in \mathcal{O}^{d\times Kp}.
\end{align}

Con la configuración anterior, ahora podemos evaluar el rendimiento en la eliminación de ruido
de un operador MSSA ligeramente modificado, que aplica un umbral a los términos que
son demasiado pequeños en la función softmax.

\begin{figure*}[t]
    \begin{center}
        \includegraphics[width=0.75\textwidth]{\toplevelprefix/chapters/chapter5/figs/MSSA.png}
        \caption{\textbf{Detalles de la arquitectura del transformador de «solo atención».}
            Cada capa consta del operador MSSA
            y una conexión de salto. Además, se incluye LayerNorm
            solo para tareas de lenguaje. En la práctica, se aplica la retropropagación
            para entrenar los parámetros del modelo utilizando muestras de entrenamiento. %
        }
        \label{fig:transformer}
    \end{center}
\end{figure*}

\begin{theorem}\label{thm:1}
    Sea $\bm Z^{1}$ tal como se define en la \Cref{def:MoG} y sea $\varphi(\cdot)$
    en \eqref{eq:MSSA}
    $\varphi(\bm x) = h\left(\sigma(\bm x)\right)$,
    donde $\sigma:\R^N \to \R^N$ es la función softmax y $h:\R^N
    \to \R^N$ es una función de umbral por elementos con $h(x) =
    \tau$ si \(x > \tau\) y \(h(x) = 0\) si \(x \leq \tau\) para
    cada $i \in [N]$.
    Supongamos que $p \gtrsim \log N$, $\delta \lesssim \sqrt{\log
    N}/\sqrt{p}$ y
    \begin{align*}
        \tau \in \left( \frac{1}{2},  \frac{1}{1+N\exp(-9p/32)} \right].
    \end{align*}
    Para valores suficientemente grandes de $N$, se cumple con una probabilidad de al menos
    $1-KLN^{-\Omega(1)}$ que, para cada $\ell \in [L]$,
    \begin{align}\label{eq:SNR}
        \mathrm{SNR}(\bm Z_k^{\ell+1}) = (1+\eta\tau)
        \mathrm{SNR}(\bm Z_k^{\ell}),\ \forall k \in [K].
    \end{align}
\end{theorem}
Este teorema demuestra que, cuando las representaciones iniciales de los tokens
se obtienen a partir de una mezcla de distribuciones gaussianas de bajo rango con un
nivel de ruido de $O(\sqrt{\log N}/\sqrt{p})$, cada
aplicación del operador MSSA modificado elimina el ruido de las representaciones de los tokens
a una tasa lineal. Cabe destacar que nuestros resultados teóricos están bien respaldados por
las observaciones experimentales realizadas con el operador MSSA \textit{real} en la
\Cref{fig:MSSA}. Este teorema proporciona una base teórica para
la capacidad práctica de eliminación de ruido de la arquitectura del transformador
derivada mediante el desenrollado de \eqref{eq:MSSA}.

\paragraph{Transformadores de solo atención.} Para verificar que esta
capacidad de eliminación de ruido se traduce en un rendimiento razonable con datos y tareas reales,
proponemos formalmente una arquitectura de transformador basada exclusivamente en la atención,
a la que denominamos AoT. Concretamente, al desenrollar
los pasos de optimización iterativa \eqref{eq:MSSA} como capas de una red profunda,
construimos una arquitectura de transformador tal y como se muestra en la
\Cref{fig:transformer}. Cada capa de la arquitectura propuesta
consta únicamente del operador MSSA y una conexión de salto. Para las tareas
lingüísticas, incorporamos además LayerNorm antes del operador MSSA
para mejorar el rendimiento. La arquitectura completa se construye
apilando dichas capas, junto con pasos esenciales de preprocesamiento
y posprocesamiento específicos de la tarea, tales como la codificación
posicional, la incrustación de tokens y una cabeza final específica de la tarea para adaptarse a
diferentes aplicaciones. La \Cref{tab:aot_vision} muestra el rendimiento de 
este transformador minimalista en la clasificación
de imágenes; el rendimiento es razonable, a pesar de su diseño
minimalista. A continuación, analizaremos una mejora conceptual de
la formulación CRATE que \textit{también} mejora el rendimiento y
la eficiencia.

\begin{table}[!t]
    \caption{\textbf{Evaluación del AoT en tareas de visión}: precisión top-1
    en ImageNet.}
    \centering
    \begin{tabular}{lcccc}
        \toprule
        Models & Accuracy  & \# of Parameters \\
        \midrule
        \midrule
        AoT   & 71.7\% & 22M \\
        CRATE & 79.5\%  & 39M  \\
        ViT & 72.4 \% & 22M \\
        \bottomrule
    \end{tabular}
    \label{tab:aot_vision}
\end{table}

\subsection{Atención en tiempo lineal: transformador de estadísticas de tokens}\label{sub:tost}

En esta subsección, proponemos un nuevo operador de atención de transformador
cuya complejidad computacional varía linealmente con el número de
tokens, basándonos en el objetivo de reducción de la tasa de codificación.  En concreto,
derivamos una nueva forma variacional del objetivo MCR$^2$ y demostramos
que la arquitectura resultante del descenso de gradiente desenrollado de
este objetivo variacional da lugar a un nuevo módulo de atención denominado
Token Statistics Self-Attention (\text{TSSA}). \text{TSSA} tiene
{\em una complejidad computacional y de memoria lineal} y se aleja radicalmente
de la arquitectura de atención típica que calcula
las similitudes por pares entre tokens.  Al sustituir el operador de atención tradicional
de un transformador por TSSA se obtiene una nueva arquitectura de red neuronal
denominada Token Statistics Transformer (ToST),
que es muy eficiente y presenta un rendimiento empírico prometedor. Como
preludio, recordemos de la ecuación \eqref{eq:MCR pi} que $\bm \Pi = [\bm
\pi_1, \ldots, \bm \pi_K] \in \R^{N \times K}$ denota una matriz estocástica
de «asignación de grupos» (es decir, $\bm \Pi \bm 1 = \bm 1$ y
$\Pi_{ik} \geq 0, \  \forall (i,k) \in [N] \times [K]$), donde
$\Pi_{ik}$ denota la probabilidad de asignar el $i$-ésimo token al
$k$-ésimo grupo.

\paragraph{Una nueva forma variacional para las tasas de codificación.} Para empezar,
consideramos una forma general de objetivos del tipo MCR$^2$ basados en funciones cóncavas
del espectro de una matriz.  Concretamente, para una matriz PSD
dada $\bm M \in \PSD(d)$ y cualquier escalar $c \geq 0$, tenemos que
$\log \det (\I + c \bm M) = \sum_{i=1}^{d} \log( 1 + c \lambda_i(\bm
M))$, donde $\lambda_i (\bm M)$ es el $i$-ésimo mayor valor propio de
$\bm M$. Además, cabe señalar que $\log(1 + c \sigma)$ es una función cóncava
y no decreciente de $\sigma$. Por lo tanto, describimos nuestros resultados
en términos de una forma más general de MCR$^2$ basada en funciones espectrales generales
de matrices PSD de la forma $F(\bm M) = \sum_{i=1}^{d}
f(\lambda_i(\bm M))$, donde $f$ es cóncava y no decreciente.  En
particular, recordemos de lo anterior en este capítulo que el mecanismo de atención
surge del desarrollo del componente de compresión de MCR$^2$,
por lo que consideramos una función de compresión más general al estilo de MCR$^2$:
\begin{equation}
    R_{c,f}(\bm Z, \bm \Pi) \doteq \frac{1}{2}\sum_{k=1}^K
    \frac{N_k}{N} F\left( \frac{1}{N_k}\Z \mathrm{Diag}(\bm \pi_k)
    \Z^{\top} \right).
    \label{eq:mcr_c_gen}
\end{equation}

Con respecto al objetivo anterior, observamos ahora el siguiente resultado:
\begin{theorem}
    \label{thm:var_concave}
    Sea \(f \colon [0, \infty) \to \R\) una función no decreciente y cóncava,
    que cumple \(f(0) = 0\) y sea \(F \colon \PSD(d) \to \R\) una función
    la forma \(F(\bm M) = \sum_{i = 1}^{d}f(\lambda_{i}(\bm M))\).
    Entonces, para cada \(\bm M \in \PSD(d)\) y \(\bm Q \in \O(d)\), tenemos
    \begin{equation}
        \label{eq:U_bound}
        F(\bm M) \leq  \sum_{i=1}^{d} f\left( (\bm Q^{\top} \bm M \bm
        Q)_{ii} \right).
    \end{equation}
    Además, la desigualdad de \eqref{eq:U_bound} se cumple con
    igualdad para cualquier $\bm Q$ que diagonalice $\bm M$, y si $f$
    es estrictamente cóncava, entonces la desigualdad de \eqref{eq:U_bound} se
    cumple con igualdad \textit{si y solo si} $\bm Q$ diagonaliza $\bm M$.
\end{theorem}

A partir del resultado anterior, podemos sustituir \eqref{eq:mcr_c_gen} por un
objetivo variacional equivalente con la forma
\begin{equation}
    \label{eq:var_comp}
    R^{\rm var}_{c,f} (\Z,\bm \Pi \mid \bm U_{[K]}) \doteq
    \frac{1}{2}\sum_{k=1}^K \frac{N_k}{N} \sum_{i=1}^{d}
    f\left(\frac{1}{N_k} (\bm U_{k}^{\top} \Z \mathrm{Diag}(\bm
    \pi_k) \bm Z^{\top} \bm U_{k})_{ii} \right),
\end{equation}
donde la equivalencia se entiende en el sentido de que, para una elección óptima de
$\{ \bm U_k \in \O(d)\}_{k=1}^K$ matrices, tal como se describe en el
\Cref{thm:var_concave} (es decir, matrices ortogonales que diagonalizan
cada $\Z \mathrm{Diag}(\bm \pi_k) \Z^\top $) lograremos un límite
estrechamente relacionado con $ R^{\rm var}_{c,f} (\Z,\bm \Pi \mid \bm U_{[K]}) =
R_{c,f} (\Z,\bm \Pi)$. Tenga en cuenta que, en general, alcanzar este límite
requeriría seleccionar, para cada instancia muestreada de $\Z$, un nuevo
conjunto óptimo de matrices de parámetros $\bm U_{k}$ que diagonalicen cada
$\Z\mathrm{Diag}(\bm \pi_{k})\Z^{\top}$, lo cual es claramente inviable
para la arquitectura de la red.
En cambio, desde un punto de vista alternativo, en lugar de considerar los
datos ($\Z$) como fijos e intentar optimizar los parámetros $\bm U_k$
para alcanzar el límite variacional más estricto, podemos adoptar el
principio de diseño de «desenrollamiento» algorítmico descrito anteriormente y diseñar un
operador que perturbe $\Z$ para minimizar de forma incremental $R_{c, f}^{\rm
var}(\cdot \mid \bm U_{[K]})$. Para dejarlo claro, cada
límite variacional se vuelve estricto cuando los espacios propios de $\Z
\mathrm{Diag}(\bm \pi_k) \Z
^\top$ se alinean con las columnas de $\bm U_k$, por lo que al rotar las
columnas apropiadas de $\Z$ (es decir, aquellas que corresponden a entradas grandes
en $\bm \pi_k$) para alinearlas con $\bm U_k$ podemos acercarnos a un
límite variacional ajustado. Es decir, en lugar de girar $\bm U_k$ para
alinearlo con los datos de cada instancia de $\Z$, podemos girar
las características de los tokens de cada $\Z$ para alinearlas con $\bm U_k$.

Siguiendo este enfoque, calculamos un paso de descenso por gradiente en
$R_{c,f}^{\rm var}$ con respecto a $\Z$.
Para iniciar este cálculo, supongamos primero que $\bm \pi \in \Re^N$ es un
vector no negativo elemento por elemento. Entonces tenemos
\begin{equation}\label{eq1:grad Rcf}
    \nabla_{\Z} \ \frac{1}{2} \sum_{i = 1}^{d} f(  (\Z
    \mathrm{Diag}(\bm \pi) \Z^{\top} )_{ii}) %
    = \; \mathrm{Diag}(\nabla f[ \Z^{\hada 2} \bm \pi] ) \Z
    \mathrm{Diag}(\bm \pi),
\end{equation}
donde $\nabla f$ es el gradiente de $f$ y (recordemos que) $\nabla
f[\cdot]$ aplica $\nabla f$ a cada elemento del vector entre
paréntesis. En particular, para $ f(x) = \log(1 + (d/\epsilon^{2}) x)$,
$\nabla f(x) = (d / \epsilon^{2}) (1+ (d / \epsilon^{2}) x)^{-1}$ es
simplemente una activación no lineal. Además, (recordemos que) $N_{k} = \langle \bm
\pi_{k}, \bm 1\rangle$. Por lo tanto, el gradiente de $R^{\rm var}_{c,f}$
respecto a $\Z$ es:
\begin{align}\label{eq2:grad Rcf}
    &\nabla_{\Z} R^{\rm var}_{c,f}(\Z,\bm \Pi \mid \bm
    U_{[K]}) \\
    &=
    \frac{1}{n}
    \sum_{k=1}^K \bm U_{k} \underbrace{\mathrm{Diag} \left(
            \nabla f\left[ (\bm
                U_{k}^{\top} \Z)^{\hada 2}  \frac{\bm \pi_k}{\langle
                    \bm \pi_{k},
    \bm 1\rangle} \right] \right)}_{\doteq \bm D(\Z, \bm \pi_{k} \mid
    \bm U_{k})} \bm U_{k}^{\top} \Z \mathrm{Diag}(\bm \pi_k).
\end{align}
(Obsérvese que la constante $1/N$ se deriva de una constante $(N_{k}/N)\cdot (1/N_{k})
= 1/N$ presente en cada término de la suma.) Si ahora consideramos un
paso de gradiente con respecto al token $j$-ésimo $\z_j$, 
llegamos a nuestro operador de compresión incremental propuesto, es decir,
nuestro sustituto de un operador de \textit{autoatención} + residual:
\begin{align}\label{eq:layer-operation-orig}
    \z_j^+
    &= \z_{j} - \tau \nabla_{\z_{j}}R_{c, f}^{\rm var}(\Z, \bm
    \Pi \mid \bm U_{[K]}) \\
    &= \z_j - \frac{\tau}{N} \sum_{k=1}^K
    \Pi_{jk} \bm U_{k} \bm D(\Z, \bm \pi_{k} \mid \bm U_{k}) \bm
    U_{k}^{\top} \z_j
\end{align}
para cada \(j \in [n]\), donde $\tau > 0$ es un parámetro de tamaño de paso para
la optimización incremental.

\paragraph{Interpretación del modelo.} Dado el operador de atención
propuesto en \eqref{eq:layer-operation-orig}, recordemos primero que las
filas de $\bm\Pi$ son no negativas y suman 1, por lo que nuestro operador 
toma un promedio ponderado de $K$ operadores similares a «cabezas de atención»
y luego añade una conexión residual. Utilizando
que \(\sum_{k = 1}^{K}\Pi_{jk} = 1\), podemos reescribir
\eqref{eq:layer-operation-orig} como: %
\begin{equation}
    \label{eq:z_up}
    \z_j^+ = \sum_{k=1}^K \Pi_{jk} \Big[ \z_j \underbrace{-
            \frac{\tau}{n} \bm U_{k} \bm D(\Z, \bm \pi_{k} \mid \bm U_{k})
    \bm U_{k}^{\top}}_{\text{action of one attention head}} \z_j \Big].
\end{equation}
Es decir, podemos considerar que cada cabeza de atención primero proyecta las características de los tokens
sobre la base $\bm U_{k}$ multiplicándolas por $\bm U_k^\top$, las multiplica por la
matriz diagonal $\bm D(\Z, \bm \pi_{k} \mid \bm U_{k})$ (abreviada como \(\bm
D_{k}\)), proyectándolas de nuevo en la base estándar mediante la multiplicación por $\bm
U_{k}$ y restando esto de las características originales del token a través de la conexión residual.
El aspecto central de nuestra capa de atención es el cálculo de $\bm
D_{k}$.  Es decir, dado que \(\Pi_{jk} \geq 0\), entonces $\bm \pi_k / \langle \bm \pi_{k}, \bm
1\rangle \in \Re^N$ forma una distribución de probabilidad que indica qué tokens pertenecen al
$k$-ésimo grupo. Como resultado, $(\bm U^\top_k \Z)^{\hada 2}
\bm \pi_k / \langle \bm \pi_{k}, \bm 1\rangle$ estima el segundo
momento de $\bm U_k^\top \Z$ bajo la distribución dada por $\bm 
\pi_k /  \langle \bm \pi_{k}, \bm 1\rangle$.  Además, dado que $f$ es una
función cóncava no decreciente, $\nabla f(x)$ disminuye monótonamente
hacia $0$ a medida que $x$ aumenta, por lo que las entradas de $\bm D_{k}$
(que tienen la forma $\nabla f(x)$) alcanzan su máximo en $x=0$ %
y decaen monótonamente a $0$ a medida que $x$ aumenta.

A partir de esto, llegamos a la interpretación fundamental de nuestros operadores de atención + residual
$[\I - (\tau/n)\bm U_{k}\bm D_{k}\bm U_{k}^{\top}]$. Es decir, este
operador realiza una proyección aproximada de rango bajo dependiente de los datos, en la que
se conservan las direcciones que tienen una gran «potencia» tras la proyección $\bm
U_k^\top \Z$ (es decir, las direcciones que tienen un segundo momento elevado $(\bm
    U_{k}^{\top}\Z)^{\hada 2}\bm \pi_k / \langle \bm \pi_{k}, \bm
1\rangle$) se conservan, mientras que las direcciones que no la tienen se
suprimen. Para comprenderlo, recordemos que los elementos de $\bm D_k$
disminuyen monótonamente hasta 0 a medida que aumenta el momento de segundo orden, por lo que, para
las direcciones con momentos de segundo orden grandes, el operador de atención + residual
actúa en gran medida como el operador de identidad. Por el contrario, para
las direcciones con un momento de segundo orden pequeño, nuestro operador resta una
proyección de los tokens a lo largo de esas direcciones, lo que da como resultado la
supresión de dichas direcciones. En comparación con el operador de autoatención
estándar, nuestro método claramente no calcula ninguna similitud
por pares entre tokens. Más bien, las interacciones entre los
elementos de $\Z$ influyen en el operador únicamente a través de su contribución
a la estadística del segundo momento utilizada para construir los $\bm D_{k}$.
No obstante, al igual que el operador de autoatención estándar, nuestro
método sigue teniendo una interpretación clara como una forma de
compresión hacia una estructura de bajo rango dependiente de los datos, en el sentido
de que realiza una proyección aproximada de bajo rango, en la que las
direcciones específicas que se suprimen son aquellas que no están
fuertemente alineadas con otros tokens del grupo.

\paragraph{Consideraciones computacionales.} Tras haber presentado nuestro
operador de atención propuesto, analizaremos ahora cómo se puede calcular
en la práctica. En primer lugar, hasta este momento de la presentación, hemos
evitado abordar cómo se «agrupan» los tokens en los distintos
cabezales de atención mediante la matriz $\bm\Pi$, pero es evidente que se necesita una forma de
construir $\bm\Pi$ para implementar nuestro método.
Además, nuestra forma variacional en el \Cref{thm:var_concave} requiere
que las matrices $\bm U$ sean cuadradas y ortogonales, pero lo ideal
sería utilizar matrices más pequeñas (es decir, reducir el número de
columnas en $\bm U$) para mejorar la eficiencia, así como eliminar las restricciones
de ortogonalidad.

En la práctica, no aplicamos las restricciones de ortogonalidad. Para
reducir el número de columnas en las matrices $\bm U$, observamos que,
al igual que en CRATE \citep{yu2023white}, si suponemos que las características $\bm
Z$ dentro del grupo \(k\) se agrupan (aproximadamente) en torno a un
subespacio de baja dimensión —digamos, de dimensión \(p\) --- entonces la
covarianza dentro del grupo \(k\) \(\Z\mathrm{Diag}(\bm
\pi_{k})\Z^{\top}\) es de rango bajo, donde recordemos que 
\cite{yu2020learning} demuestra que la geometría óptima de $\Z$ se da 
cuando cada grupo es un subespacio de rango bajo, ortogonal a los demás
grupos. Por lo tanto, podemos encontrar explícitamente una base ortonormal de baja dimensión
para la imagen de esta covarianza, es decir, el espacio lineal generado por los
datos del grupo \(k\). Si la dimensión es $p \leq d$, la base puede
representarse mediante una matriz ortogonal $d\times p$ $\bm U_k \in \O(d,
p)$. En este caso, podemos establecer una cota superior más eficiente para \(R_{c,f}\)
utilizando estas matrices de base ortogonal de rango bajo. Para demostrarlo, utilizamos
una versión más general del \Cref{thm:var_concave}  para obtener el
siguiente corolario.
\begin{corollary}\label{cor:var_concave_logdet}
    Sea \(f \colon [0, \infty) \to \R\) una función no decreciente y cóncava,
    que cumple \(f(0) = 0\) y sea \(F \colon \PSD(p) \to \R\) una función
    la forma \(F(\bm M) = \sum_{i = 1}^{p}f(\lambda_{i}(\bm M))\).
    Sean \(\Z\), \(\bm \Pi\) fijos. Entonces, para todo \(\bm U_{1},
    \dots, \bm U_{K} \in \O(d, p)\) tal que
    \(\mathrm{image}(\Z\diag(\bm \pi_{k})\Z^{\top}) \subset
    \mathrm{image}(\bm U_{k})\) para todo \(k \in [K]\), tenemos
    \begin{equation}
        R_{c, f}(\Z, \bm\Pi) \leq R_{c, f}^{\rm var}(\Z, \bm\Pi \mid
        \bm U_{[K]}),
    \end{equation}
    donde \(R_{c, f}^{\rm var}\) se define formalmente en
    \eqref{eq:var_comp}. La igualdad se cumple si \(\bm U_{k}\) diagonaliza
    \(\Z\diag(\bm \pi_{k})\Z^{\top}\) para cada \(k \in [K]\) y si
    \(f\) es fuertemente cóncava, entonces esta condición de igualdad se convierte en un
    ``si y solo si.''
\end{corollary}

El último paso para definir nuestro operador de atención consiste en estimar la
pertenencia al grupo $\bm\Pi$. Para ello, planteamos un modelo sencillo de cómo
cada característica \(\z_{j}\) se desvía de su subespacio de soporte y
luego hallamos la asignación óptima de subespacios. \cite{yu2023white} muestran
que si modelamos independientemente cada \(\z_{j}\) como perteneciente a un
modelo de mezcla gaussiana de baja dimensión, donde cada gaussiana tiene una
matriz de covarianza con espectro idéntico y está apoyada en un
subespacio con base ortonormal \(\bm U_{k}\), más
ruido gaussiano independiente con covarianza \(\eta \I\), entonces la probabilidad a posteriori
de que cada token \(\z_{j}\) pertenezca a cada subespacio viene
dada por la matriz de asignación \(\bm \Pi = \bm \Pi(\bm Z \mid \bm
U_{[K]})\) de la siguiente manera:
\begin{align}\label{eq:pi}
    \bm \Pi & =
    \begin{bmatrix} \boldsymbol{\nu}(\z_1 \mid \bm U_{[K]})^\top
        \\ \vdots \\ \boldsymbol{\nu}(\z_n \mid \bm U_{[K]})^\top
    \end{bmatrix}, \quad
    \text{where} \\
    \boldsymbol{\nu}(\z_j \mid \bm U_{[K]}) &\doteq
    \operatorname{softmax}\left( \frac{1}{2\eta}
        \begin{bmatrix} \|  \bm U_{1}^{\top} \z_j \|_2^2 \\ \vdots
            \\ \| \bm U_{K}^{\top} \z_j \|_2^2
    \end{bmatrix} \right), \quad \forall j \in [n],
\end{align}
donde $\eta$ se convierte en un parámetro de temperatura aprendible. Así, dada
una característica de entrada \(\Z\), estimamos \(\bm\Pi\) utilizando \eqref{eq:pi}
y luego calculamos el operador de atención. Al combinar la construcción
de $\bm\Pi$ en \eqref{eq:pi} con
\eqref{eq:layer-operation-orig}, obtenemos el operador {\em Token Statistics
Self-Attention}:
\begin{equation}
    \label{eq:tssa}
    \operatorname{TSSA}(\Z \mid \bm U_{[K]}) \doteq -\frac{\tau}{n}\sum_{k
    = 1}^{K}\bm U_{k}\bm D(\Z, \bm\pi_{k} \mid \bm U_{k})\bm
    U_{k}^{\top}\Z\diag(\bm \pi_{k}),
\end{equation}
donde \(\bm\pi_{k}\) son las columnas de \(\bm\Pi = \bm\Pi(\Z \mid \bm
U_{[K]})\) definidas en \eqref{eq:pi} y \(\bm D\) se define en
\eqref{eq2:grad Rcf}.

Partiendo de una arquitectura Transformer estándar y sustituyendo
el operador de atención por TSSA, obtenemos una arquitectura denominada
Token Statistics Transformer (ToST), que se muestra en la
\Cref{fig:vcrate-architecture}. Nos gustaría volver a destacar que
el operador TSSA reduce la complejidad temporal y de memoria habitual del
operador de atención, en particular al reducir la complejidad temporal y
de memoria cuadrática a lineal, lo que hace que ToST sea mucho más eficiente
que su homólogo diseñado empíricamente. A continuación, mostramos brevemente el
rendimiento y la eficiencia empíricos de la arquitectura ToST.

\begin{figure}[!t]
    \centering
    \includegraphics[width=1\textwidth,trim={0 5.0cm 0
    0},clip]{\toplevelprefix/chapters/chapter5/figs/V-CRATE.pdf}
    \vspace{-1mm}
    \caption{\small \textbf{Una capa $\ell$ del Token
        Statistics Transformer (ToST) propuesto.} Cabe destacar que la autoatención
        de ToST transforma los tokens $\Z^{\ell}$ de manera eficiente en
        $\Z^{\ell+1}$, multiplicando cada fila del token proyectado
        por \textit{solo un escalar}. Esto conduce a una
        complejidad reducida de la atención: tiene una complejidad espacial de $O(p)$ y una complejidad temporal de $O(pn)$,
        donde $p$ es la dimensión de los tokens proyectados
        de cada cabeza y $n$ es el número de tokens.
    }
    \label{fig:vcrate-architecture}
\end{figure}

\begin{figure}[!t]
    \centering
    \begin{subfigure}{0.475\textwidth}
        \includegraphics[width=0.49\textwidth]{\toplevelprefix/chapters/chapter5/figs/tost_inference_time_vit.pdf}
        \includegraphics[width=0.49\textwidth]{\toplevelprefix/chapters/chapter5/figs/tost_gpu_memory_vit.pdf}
        \caption{}
    \end{subfigure}
    \hfill
    \begin{subfigure}{0.475\textwidth}
        \includegraphics[width=0.49\textwidth]{\toplevelprefix/chapters/chapter5/figs/tost_inference_time_gpt2.pdf}
        \includegraphics[width=0.49\textwidth]{\toplevelprefix/chapters/chapter5/figs/tost_gpu_memory_gpt2.pdf}
        \caption{}
    \end{subfigure}
    \caption{\small\textbf{Complejidad temporal y de memoria de ToST en
        la práctica} en comparación con ViT (a) y GPT-2
        (b), diseñados empíricamente. Observamos que ToST solo requiere tiempo y
    memoria lineales, en comparación con los costos cuadráticos de los transformadores tradicionales.}
\end{figure}

\begin{table*}[t!]
    \centering
    \caption{\small \textbf{Precisión «top-1» de ToST en diversos
            conjuntos de datos con diferentes tamaños de modelo tras el preentrenamiento en
        ImageNet.} Para ImageNet/ImageNetReaL, evaluamos directamente la
        precisión «top-1». Para otros conjuntos de datos, preentrenamos el modelo en
        ImageNet y luego evaluamos el rendimiento del aprendizaje por transferencia mediante
        el ajuste fino. Observamos que el rendimiento de ToST escala con el
    tamaño del modelo y es competitivo con el ViT diseñado empíricamente.}
    \label{tab:tost_in1k_transfer}
    \small
    \setlength{\tabcolsep}{10.6pt}
    \resizebox{0.95\textwidth}{!}{%
        \begin{tabular}{@{}lccc|cc@{}}
            \toprule
            Datasets & ToST-T(iny) & ToST-S(mall) &
            ToST-M(edium) &  ViT-S &
            ViT-B(ase)\\
            \toprule
            \# parameters & 5.8M & 22.6M & 68.1M &  22.1M & 86.6 M  \\
            \midrule
            \midrule
            ImageNet & 67.3 & 77.9 & 80.3 & 79.8 & 81.8 \\
            ImageNet ReaL & 72.2  & 84.1 & 85.6 &  85.6 & 86.7 \\
            \midrule
            CIFAR10 & 95.5 & 96.5 & 97.5 &  98.6 & 98.8 \\
            CIFAR100 & 78.3 & 82.7 & 84.5 &  88.8 & 89.3 \\
            Oxford Flowers-102 & 88.6 & 92.8 & 94.2 &  94.0 & 95.7 \\
            Oxford-IIIT-Pets & 85.6 & 91.1 & 92.8 & 92.8 & 94.1 \\
            \bottomrule
        \end{tabular}%
    }
\end{table*}

\begin{table*}
    \centering
    \caption{\small\textbf{Rendimiento de ToST en
        tareas de contexto largo.} Demostramos el rendimiento en Long Range Arena, un
        banco de pruebas para tareas de contexto largo (\(n\) grande). En este banco de pruebas,
        ToST presenta un rendimiento y una eficiencia claramente superiores en comparación con
    todas las variantes de Transformer, incluido el propio Transformer.}
    \label{tab:tost_lra}
    \resizebox{0.8\textwidth}{!}{%
        \begin{tabular}{@{}lccccccc@{}}
            \toprule
            Model        & ListOps  & Text     & Retrieval & Image
            & Pathfinder & Avg      \\
            \midrule
            \midrule
            Reformer              & 37.27 & 56.10
            & 53.40              & 38.07             & 68.50
            & 50.56             \\
            BigBird               & 36.05             & 64.02
            & 59.29              & 40.83             & 74.87
            & 54.17             \\
            LinFormer         & 16.13             & 65.90
            & 53.09              & 42.34             &
            75.30               & 50.46             \\
            Performer             & 18.01             & 65.40
            & 53.82              & 42.77             & 77.05
            & 51.18             \\
            Transformer           & 37.11             & 65.21
            & 79.14              & 42.94
            & 71.83              & 59.24            \\
            ToST & 37.25    &
            66.75    & 79.46     & 46.62
            &    69.41      &     59.90 \\
            \bottomrule
        \end{tabular}%
    }
\end{table*}

\subsection{CRATE causal}
\label{sec:causal-crate}
Algunas de las aplicaciones más populares de las redes profundas se dan en el
ámbito de los \textit{datos secuenciales}, donde, a diferencia del ejemplo anterior
de las imágenes, los datos siguen un orden natural. Por ejemplo, a diferencia
de las imágenes (donde se tiende a observar la imagen completa de una sola vez),
el lenguaje se procesa token por token (de manera causal). Este simple
hecho impulsa el desarrollo de modelos de lenguaje a gran escala (LLM), que
se entrenan para predecir (la distribución de probabilidad del) siguiente
token de una secuencia basándose \textit{únicamente en el historial de los tokens 
anteriores}. A continuación, para obtener una muestra de un LLM, basta con realizar un muestreo iterativo
a partir de la distribución prevista para el siguiente token y, a continuación, añadir
el token muestreado al historial. Para entrenar LLM y otros
modelos, como los de generación de vídeo, es necesario desarrollar
arquitecturas capaces de realizar este cálculo causal de manera eficiente.

Teniendo esto en cuenta y basándonos en la comprensión de los procesos autorregresivos causales
del \Cref{ch:intro}, decimos que un codificador \(f\) es
\textit{causal} si, para cada entrada \(\vX = [\vx_{1}, \dots, \vx_{N}]\), se cumple que
\begin{align}
    f(\vX)_{1:N-1}
    &= f([\vx_{1}, \dots, \vx_{N}])_{1:N-1} =
    f([\vx_{1}, \dots, \vx_{N - 1}]) \\
    &= f(\vX_{1:N-1}),
\end{align}
donde \(\vZ_{1:N-1} = [\vz_{1}, \dots, \vz_{N-1}]\), etc. En español,
esto significa que las primeras \(N-1\) características son iguales a la salida del
codificador sobre las primeras \(N-1\) incrustaciones de tokens de la entrada; dicho de manera aún
más sencilla, significa que podemos calcular las características de cada token
\textit{una por una} y esto sería equivalente a calcular las
características de toda la entrada. Esta causalidad es necesaria para las
aplicaciones mencionadas anteriormente (por ejemplo, los LLM). En particular, las
implementaciones de CRATE y ToST presentadas anteriormente
\textit{no son causales} (la demostración se deja como ejercicio).

Podemos construir arquitecturas causales de white-box mediante diversos
métodos. Aquí presentaremos un método sencillo que se basa en nuestro
marco de optimización desenrollada anterior. Específicamente, calculamos
las características \(\vz_{1}, \dots, \vz_{N}\) correspondientes a los primeros
\(N\) tokens de la entrada \textit{uno por uno} para optimizar el
objetivo de aprendizaje de la representación, como la reducción dispersa de la tasa 
\eqref{eq:sparse-rr-1}:
\begin{equation}
    \max_{\vz_{i}}\ \Delta R(\vZ_{1:i} \mid \vU_{[K]}) - \lambda
    \|\vZ_{1:i}\|_{1},\ \forall i \in [N].
\end{equation}
Si seguimos el procedimiento de desenrollamiento en dos pasos que
dio lugar a CRATE, podemos obtener la iteración:
\begin{align}
    \vz_{i}^{\ell + 1/2} &\approx \vz_{i}^{\ell} - \kappa
    \nabla_{\vz_{i}} R(\vZ_{1:i}^{\ell} \mid \vU_{[K]}^{\ell}), \\
    \vz_{i}^{\ell + 1} &\approx \argmin_{\vz_{i}} \left\{\lambda
        \|\vz_{i}\|_{1} + \frac{1}{2}\norm{\vz_{i}^{\ell + 1/2} -
    \vD^{\ell} {\vz_{i}}}_F^2\right\},
\end{align}
o bien, utilizando la misma conversión de estos dos pasos a operadores de red
que usamos para CRATE,
\begin{align}
    \vz_{i}^{\ell + 1/2}
    &= \left(1 - \frac{\kappa p}{N\epsilon^{2}}\right)\vz_{i}^{\ell}
    + \frac{\kappa p}{N
    \epsilon^{2}}\operatorname{MSSA}(\vZ_{1:i}^{\ell} \mid
    \vU_{[K]}^{\ell})_{i}, \\
    \vz_{i}^{\ell + 1}
    &= \operatorname{ISTA}(\vz_{i}^{\ell + 1/2} \mid \vD^{\ell}).
\end{align}
Analicemos esta iteración con un poco más de detalle. En primer lugar,
observemos que, por definición, esta secuencia de características corresponde
a un codificador causal. A continuación, supongamos que nos encontramos en un escenario
en el que hemos calculado \(\vZ_{1:i-1}^{\ell + 1} = [\vz_{1}^{\ell +
1}, \dots, \vz_{i-1}^{\ell + 1}]\), habiendo calculado las cantidades
\begin{equation}
    (\vU_{k}^{\ell})^{\top}\vZ_{1:i-1}^{\ell} =
    [(\vU_{k}^{\ell})^{\top}\vz_{1}^{\ell}, \dots,
    (\vU_{k}^{\ell})^{\top}\vz_{i-1}^{\ell}]
\end{equation}
a lo largo del proceso y queremos calcular \(\vz_{i}^{\ell + 1}\) (por
ejemplo, en el caso de la inferencia de un modelo de lenguaje grande). En este caso, cabe señalar
que la regla de actualización para \(\vz_{i}^{\ell + 1/2}\) puede simplificarse de la siguiente manera:
\begin{align}
    &\vz_{i}^{\ell + 1/2} \\
    &= \left(1 - \frac{\kappa p}{N\epsilon^{2}}\right)\vz_{i}^{\ell}
    + \frac{\kappa p}{N
    \epsilon^{2}}\operatorname{MSSA}(\vZ_{1:i}^{\ell} \mid
    \vU_{[K]}^{\ell})_{i} \\
    &= \left(1 - \frac{\kappa p}{N\epsilon^{2}}\right)\vz_{i}^{\ell}
    + \frac{\kappa p^{2}}{N^{2} \epsilon^{4}}\sum_{k =
    1}^{K}\vU_{k}\vU_{k}^{\top}\vZ_{1:i}\softmax((\vU_{k}^{\top}\vZ_{1:i})^{\top}(\vU_{k}^{\top}\vz_{i}))
    \\
    &= \scalebox{0.84}{\(\displaystyle\left(1 - \frac{\kappa
            p}{N\epsilon^{2}}\right)\vz_{i}^{\ell}
            + \frac{\kappa p^{2}}{N^{2} \epsilon^{4}}\sum_{k =
            1}^{K}\vU_{k}\vU_{k}^{\top}[\vZ_{1:i-1},
            \vz_{i}]\softmax((\vU_{k}^{\top}[\vZ_{1:i-1},
    \vz_{i}])^{\top}(\vU_{k}^{\top}\vz_{i}))\)} \\
    &= \scalebox{0.7625}{\(\displaystyle\left(1 - \frac{\kappa
            p}{N\epsilon^{2}}\right)\vz_{i}^{\ell}
            + \frac{\kappa p^{2}}{N^{2} \epsilon^{4}}\sum_{k =
            1}^{K}\vU_{k}[\vU_{k}^{\top}\vZ_{1:i-1},
            \vU_{k}^{\top}\vz_{i}]\softmax(([\vU_{k}^{\top}\vZ_{1:i-1},
    \vU_{k}^{\top}\vz_{i}])^{\top}(\vU_{k}^{\top}\vz_{i})).\)}
\end{align}
Este paso se vuelve \textit{muy eficiente} si almacenamos en caché
\(\vU_{k}^{\top}\vZ_{1:i-1}\) de los pasos anteriores y le añadimos una sola
columna cada vez que calculamos la proyección
\(\vU_{k}^{\top}\vz_{i}\). Es decir, reducimos considerablemente el número de
productos de matrices grandes y los reemplazamos por cargas de caché y
productos de matriz-vector, lo cual es, en general, mucho más económico en términos de
complejidad temporal. Este almacenamiento en caché es la razón por la que los modelos generativos causales,
como los LLM, pueden muestrear de manera eficiente miles de tokens por
segundo, incluso si cada paso del entrenamiento tarda unos segundos por sí solo. La
caché para las proyecciones en el subespacio de las características también se conoce como
la así denominada «caché KV».

Por último, consideremos el caso en el que debemos \textit{entrenar} un
modelo CRATE causal y queremos encontrar la forma más eficiente de hacerlo
simultáneamente, dada una secuencia de entrada completa \(\vX = [\vx_{1}, \dots, \vx_{N}]\).
El paso ISTA se puede paralelizar para convertirse en el paso
ISTA regular del CRATE no causal, es decir,
\begin{equation}
    \operatorname{ISTA}(\vZ \mid \vD) \doteq [\ISTA(\vz_{1} \mid
    \vD), \dots, \ISTA(\vz_{N} \mid \vD)],
\end{equation}
lo que implica, por tanto, que el paso ISTA sigue siendo el mismo que el
CRATE no causal. El paso MSSA resulta más interesante, ya que cambia
de manera significativa. Para ver cómo cambia, cabe señalar que cada operador MSSA
tiene una estructura interesante que merece cierta atención.
Concretamente, si definimos el \textit{operador MSSA causal} como la matriz en bloque:
\begin{equation}
    \operatorname{CMSSA}(\vZ \mid \vU_{[K]}) \doteq
    [\operatorname{MSSA}(\vZ_{1:1} \mid \vU_{[K]})_{1}, \dots,
    \operatorname{MSSA}(\vZ_{1:N} \mid \vU_{[K]})_{N}],
\end{equation}
entonces, analizando el álgebra de softmax (la demostración se deja, una vez más,
como ejercicio), podemos demostrar que
\begin{align}
    \operatorname{CMSSA}(\vZ \mid \vU_{[K]})
    &\doteq
    \scalebox{0.85}{\(\displaystyle\frac{p}{N\epsilon^{2}}\mat{\vU_{1}, \dots,
            \vU_{K}}\mat{\softmax((\vU_{1}^{\top}\vZ)^{\top}(\vU_{1}^{\top}\vZ)
                +\vM_{N}) \\ \vdots
    \\ \softmax((\vU_{K}^{\top}\vZ)^{\top}(\vU_{K}^{\top}\vZ) + \vM_{N})},\)} \\
    \text{where}\ &\vM_{N} = \mat{0 & -\infty & -\infty & \dots
        & -\infty & -\infty \\ 0 & 0 & -\infty & \dots & -\infty &
        -\infty \\ \vdots & \vdots & \vdots & \ddots & \vdots & \vdots
    \\ 0 & 0 & 0 & \dots & 0 & -\infty \\ 0 & 0 & 0 & \dots & 0 & 0}.
\end{align}
Aquí, la matriz \(\vM_{N} \in \R^{N \times N}\) es otra forma de
codificar que ninguna característica del CMSSA, es decir, el MSSA causal, depende de una
característica futura, ya que las entradas relevantes del argumento de la
función softmax son \(-\infty\) y por lo tanto se establecen en \(0\) después de la
función exponencial dentro de la softmax, con lo que, en la práctica, se ignoran.
Esta matriz \(\vM_{N}\) se denomina a veces \textit{máscara de atención (causal)}.

La consecuencia de esto es que, cuando disponemos de la secuencia completa
\(\vZ^{\ell}\), podemos calcular \(\vZ^{\ell + 1}\) en un tiempo similar
o inferior al de la capa CRATE habitual (ya que \(\vM_{N}\) está
    codificada de forma fija, es independiente de la entrada y nos permite ignorar muchas entradas
para la función softmax). Por lo tanto, podemos definir y entrenar el modelo
completo de secuencia a secuencia de la misma manera que la arquitectura regular.

Cuando aplicamos esto al modelado del lenguaje (los detalles se proporcionan en el
\Cref{ch:applications}), observamos que podemos obtener resultados razonables
en comparación con modelos de lenguaje de tamaño similar diseñados de forma empírica,
tal y como se muestra en la \Cref{tab:causal-crate-results}.

\begin{table}[t!]
    \small
    \caption{\small Pérdida de entropía cruzada en la validación sin entrenamiento del
        modelo Causal-CRATE-Base y de los modelos GPT2-Small y GPT2-Base,
        evaluada en la parte de prueba de los conjuntos de datos (cuanto menor, mejor).
    }
    \centering
    \resizebox{1\columnwidth}{!}
    {
        \begin{tabular}{ccccccc}
            \hline
            & \#parameters & OWT & LAMBADA &
            WikiText & PTB & Avg \\
            \hline
            \hline
            GPT2-Base  & {124M} & 2.85 & 4.12 & 3.89 & 4.63 & 3.87 \\
            {GPT2-Small } &  {64M} & {3.04} & {4.49} & {4.31} &
            {5.15} & {4.25} \\
            Causal-CRATE-Base & {60M} & 3.37 & 4.91 & 4.61 & 5.53 & 4.61 \\
            \hline
        \end{tabular}
    }
    \label{tab:causal-crate-results}
\end{table}

A pesar de las ventajas teóricas de la arquitectura causal CRATE,
la brecha en las tareas de modelado del lenguaje (y, en términos más generales, de modelado de secuencias)
sigue existiendo. A continuación, analizamos dos estrategias que podrían
cerrar esta brecha; ambas son objeto de líneas de investigación activas.

\begin{example}[Mejores objetivos y estrategias de desenrollado]
    Una posibilidad a considerar es que el objetivo de la reducción de la tasa de dispersión
    no sea el más adecuado para los datos de secuencias causales. Por ejemplo,
    podríamos querer asegurarnos de que las características tengan algún tipo de
    relación o estructura temporal, como por ejemplo, que sean generadas por un
    sistema dinámico lineal con inicialización (de mezcla) gaussiana 
    y ruido de paso gaussiano. Existe una línea de investigación reciente, pero de gran
    popularidad, dedicada a modelar de manera eficiente y explícita
    la dinámica en las características mediante los denominados \textit{modelos de espacio de estado}
    (SSM, denominados así en contravención de la noción clásica
    de «modelos de espacio de estado» discutida en el \Cref{ch:intro})
    \citep{gu2024mamba,yang2024parallelizing,yang2024gated}. Estos
    modelos son prometedores desde el punto de vista empírico porque evitan la necesidad de
    una atención entre todos los pares, eludiendo así la complejidad computacional y de memoria
    del operador MSSA y de su equivalente convencional, la
    atención multicabezal. A pesar de esta mejora en la eficiencia, son
    muy eficaces e incluso se han utilizado en algunos modelos de lenguaje a gran escala
    \citep{team2025kimi}. Mediante el diseño de objetivos más
    sofisticados y estrategias de «desenrollamiento», podríamos
    construir una arquitectura SSM de «white-box», o un híbrido SSM-transformador,
    y así reducir la brecha en las tareas causales.
\end{example}

\begin{example}[Estrategias de cálculo en tiempo de inferencia]
    Otra forma de mejorar los modelos causales consiste en utilizar más estrategias denominadas
    \textit{de cálculo en tiempo de inferencia}: durante la inferencia
    (como el muestreo a partir de un modelo de lenguaje grande, o LLM), hay que encontrar la manera de dedicar más
    recursos computacionales para obtener un mejor resultado. Si logramos
    hacerlo, podremos resolver problemas más complejos asignando más
    recursos computacionales a la tarea en cuestión.

    Un enfoque sencillo para esto consiste en aplicar un bucle al modelo: introducir las características finales
    en las capas iniciales de secuencia a secuencia y permitir
    que pase por más rondas de eliminación iterativa de ruido. Esta
    metodología se conoce en la literatura como «transformadores en bucle»
    (en el caso de que el modelo base sea un
    transformador diseñado empíricamente); se ha demostrado que este enfoque
    ofrece ventajas en algunas tareas especialmente difíciles
    \citep{giannou2023looped,yang2023looped,saunshi2025reasoning}.
    Sin embargo, desde nuestro punto de vista, resulta más razonable recorrer
    la \textit{misma capa} varias veces, aplicando
    varios pasos de eliminación de ruido o compresión locales sobre el
    \textit{mismo modelo} en cada capa antes de pasar a la siguiente.

    Otra metodología posible se basa en nuestra comprensión
    de la estructura de baja dimensión de las características y los parámetros
    del modelo. Por ejemplo, dado que se supone que las matrices \(\vU_{[K]}\)
    son soportes de los subespacios generados por las
    características, podríamos actualizarlas en el momento de la inferencia utilizando un algoritmo en línea
    como el (G)PCA en línea. Medidas similares podrían permitir
    adaptar el diccionario en el momento de la inferencia. Esto nos permitiría
    mejorar y especializar nuestros operadores locales de eliminación de ruido o compresión
    en cada capa para las características, lo que, por lo tanto, podría
    dar lugar a un mejor rendimiento.

    Se han realizado numerosos intentos empíricos para mejorar los métodos de cálculo en tiempo de inferencia.
    Este sigue siendo un campo de investigación muy activo en la actualidad.
\end{example}

% \subsection{Causal ToST} \label{sec:tost}
% \yima{Causal version of ToST: relations to nonlinear filtering, RNNs,
% and other linear time architectures...}

\section{Resumen y Notas}

Los materiales presentados en este capítulo se basan en una serie de
trabajos recientes sobre este tema, entre los que se incluyen \cite{chan2021redunet,
wang2024global, wang2025attention, wu2025token, yu2023white}. Estas
contribuciones abarcan tanto avances teóricos como metodologías prácticas
para construir redes profundas interpretables mediante
la optimización desenrollada. Muchos de los resultados y demostraciones clave que se analizan
en este capítulo se derivan directamente de estos
trabajos fundamentales o se inspiran en ellos.

La idea de «desenrollar» un algoritmo de optimización para construir una red neuronal
se remonta al trabajo pionero \cite{gregor2010learning}. En
este trabajo, los autores demostraron que los algoritmos de codificación
dispersa---como el algoritmo de umbralización y contracción iterativa
(ISTA)---pueden desenrollarse para formar perceptrones multicapa (MLP),
uniendo de manera efectiva la optimización iterativa y el diseño de redes neuronales.
Cabe destacar que \cite{monga2019algorithm} demostró que estas
redes «desenrolladas» son más interpretables, eficientes y efectivas
en comparación con las redes genéricas.

En este capítulo, partimos de esta perspectiva para desarrollar arquitecturas de redes profundas de caja abierta basadas en principios,
mediante la aplicación de algoritmos de optimización
diseñados para minimizar objetivos bien fundamentados,
como el objetivo de reducción (dispersa) de la tasa presentado
anteriormente. Este enfoque no solo aclara la función de cada capa en
la red, sino que también ofrece una base teórica para las decisiones
arquitectónicas, yendo más allá del método empírico de prueba y error hacia un diseño
interpretable y orientado a objetivos. A continuación, comparamos las redes neuronales profundas (DNN) convencionales,
que suelen construirse mediante un diseño empírico y
un ajuste heurístico, con nuestras arquitecturas ReduNet, que cuentan con una base matemática:

\begin{center}
    \begin{tabular}{| l || c | c |}
        \hline
        & DNNs convencionales & ReduNets\\ [0.5ex]
        \hline \hline
        Objetivos & Ajustes de entrada/salida & ganancia de información \\ [0.5ex]
        \hline
        Arquitecturas profundas & prueba y error & optimización iterativa \\  [0.5ex]
        \hline
        Operadores de capa & empirica & gradiente proyecteda \\  [0.5ex]
        \hline
        Invariancia al traslado & CNNs y aumento & ReduNets invariantes \\  [0.5ex]
        \hline
        Inicializaciones & aleatorio/pre-diseño & desenrollado hacia adelante \\ [0.5ex]
        \hline
        Entrenamiento/ajuste fino & retroprop & h.adelante/retroprop\\ [0.5ex]
        \hline
        Interpretabilidad & black-box & white-box \\ [0.5ex]
        \hline
        Representaciones & oculta/latente & subespacios incoherentes \\ [0.5ex]
        \hline
    \end{tabular}
\end{center}

\section{Ejercicios y extensiones}

\begin{exercise}
    Sea $\bm Z = [\bm Z_1,\dots,\bm Z_K] \in \R^{d\times N}$ tal que
    $\bm Z_k \in \R^{d\times N_k}$ para todo $k \in [K]$. Para algún
    $\alpha > 0$, sea
    \begin{align*}
        R(\bm Z) = \log\det\left(\bm I + \alpha\bm Z\bm Z^\top \right).
    \end{align*}
    \begin{enumerate}
        \item Dada una dirección cualquiera $\bm D \in \R^{d\times m}$, demuestre
            que $\nabla R(\bm Z) = \alpha\bm X^{-1}\bm Z$ y
            \begin{align*}
                &\nabla^2 R(\bm Z)[\bm D, \bm D] \\
                &=\alpha \mathrm{Tr}\left( \bm
                X^{-1}\bm D\bm D^\top\right) -
                \frac{\alpha^2}{2}\mathrm{Tr}\left(\bm X^{-1}\left( \bm Z\bm
                    D^\top+\bm D\bm Z^\top\right) \bm X^{-1}\left(
                        \bm Z\bm D^\top+\bm
                D\bm Z^\top\right)\right),
            \end{align*}
            donde $\bm X \doteq \bm I + \alpha\bm Z\bm Z^\top$. {\em
            Pista:} Obsérvese que
            \begin{align*}
                \nabla^2 R(\bm Z)[\bm D, \bm D] \doteq \left\langle \lim_{t
                \to 0} \frac{\nabla R(\bm Z+ t\bm D) - \nabla R(\bm Z)}{t},
                \bm D \right\rangle.
            \end{align*}
        \item
            Muestre que
            \begin{align*}
                R(\bm Z) \le \sum_{k=1}^K \log\det\left(\bm I + \alpha\bm
                Z_k\bm Z_k^\top \right),
            \end{align*}
            donde la igualdad se cumple si y solo si $\bm Z_k^\top\bm Z_l = \bm
            0$ para todo $k \neq l \in [K]$.
        \item Dado $\alpha > 0$; sea $\alpha_k = N\alpha/N_k$ para cada $k
            \in [K]$. Derive la fórmula cerrada para el punto crítico de primer orden
            de la siguiente función:
            \begin{align*}
                f(\bm Z_k) = \frac{1}{2}\log\det\left(\bm I + \alpha\bm
                Z_k\bm Z_k^\top \right) - \frac{N_k}{2N}\log\det\left(\bm I +
                \alpha_k \bm Z_k\bm Z_k^\top \right) -
                \frac{\lambda}{2}\|\bm Z_k\|_F^2.
            \end{align*}
            {\em Pista:} Sea $r_k=\mathrm{rank}(\bm Z_k)$. Consideremos la
            siguiente descomposición en valores singulares de $\bm Z_k$:
            \begin{equation}
                \bm Z_k = \bm P_k\bm \Sigma_k\bm Q_k^\top =
                \begin{bmatrix}
                    \bm P_{k,1} & \bm P_{k,2}
                \end{bmatrix}
                \begin{bmatrix}
                    \tilde{\bm \Sigma}_{k} & \bm 0 \\
                    \bm 0 & \bm 0
                \end{bmatrix}
                \begin{bmatrix}
                    \bm Q_{k,1}^\top \\ \bm Q_{k,2}^\top
                \end{bmatrix},
            \end{equation}
            donde $\bm P_k \in \mathcal{O}^d$ con $\bm P_{k,1} \in
            \R^{d\times r_k}$ y  $\bm P_{k,2} \in \R^{d\times (d-r_k)}$,
            $\bm \Sigma_k \in \R^{d\times N_k}$ con $\tilde{\bm \Sigma}_{k}
            \in \R^{r_k\times r_k}$ siendo una matriz diagonal, y  $\bm Q_k
            \in \mathcal{O}^{N_k}$ con $\bm Q_{k,1} \in \R^{N_k\times r_k}$
            y  $\bm P_{k,2} \in \R^{N_k \times (N_k-r_k)}$.
    \end{enumerate}
\end{exercise}

\begin{exercise}[Serie de Neumann para la inversa de una matriz]\label{ex:neumannn}
    Sea $\bm A \in \mathbb{R}^{n\times n}$. Si $\|\bm A\| < 1$, muestre que
    \begin{align}\label{eq:neumann}
        \left( \bm I - \bm A\right)^{-1} = \sum_{k=0}^\infty \bm A^k.
    \end{align}
    {\em Pista:} La demostración consta de dos pasos.
    \begin{enumerate}[(i)]
        \item {\bf Paso 1}: Demuestre que la serie infinita
            $\sum_{k=0}^\infty \bm A^k$ converge cuando $\|\bm A\| < 1$ utilizando
            $\|\bm A^k\| \le \|\bm A\|^k$.
        \item {\bf Paso 2}: Calcule el producto matricial $(\bm I - \bm A)
            \sum_{k=0}^\infty \bm A^k$.
    \end{enumerate}
\end{exercise}

% \begin{exercise}
%     In this exercise, we will consider some modern empirically
%     designed variants of attention and interpret them through a white-box lens.
% 
%     \DP{TODO @Sam: RoPE and TTT/DeltaRule}
% \end{exercise}

\begin{exercise}
    En este ejercicio analizaremos una medida de la tasa de codificación
    \(R_{\epsilon}^{c, \mu}\) que utiliza las medias de las
    distribuciones, es decir, una tasa de codificación para tokens extraídos (no i.i.d.)
    del modelo de mezcla gaussiana de rango bajo
    \begin{equation}
        \vz_{i} \sim \frac{1}{K}\sum_{k = 1}^{K}\dNorm(\vmu_{k},
        \vU_{k}\vU_{k}^{\top}).
    \end{equation}
    \begin{enumerate}
        \item Demuestre que una tasa de codificación adecuada para \(\vZ\)
            en este modelo es|
            \begin{equation}
                R_{\epsilon}^{c, \mu}(\vZ) = \sum_{k =
                1}^{K}R_{\epsilon}(\vU_{k}^{\top}(\vZ - \vmu_{k}\vone^{\top})).
            \end{equation}
        \item De manera similar a cómo derivamos el MSSA, derive un nuevo
            operador similar a la atención a partir de esta tasa de codificación. ¿Cuáles son
            las diferencias clave?
    \end{enumerate}
\end{exercise}

\begin{exercise}
    Calcule los gradientes en \eqref{eq1:grad Rcf} y
    \eqref{eq2:grad Rcf}.
\end{exercise}

\begin{exercise}
    Por favor, demuestre el \Cref{cor:var_concave_logdet} cuando $Kp \le d$.
\end{exercise}

\begin{exercise}
    Derive una versión causal del operador TSSA.
\end{exercise}

\end{document}
